\documentclass{article}

\usepackage{url}
\usepackage{amssymb}
\usepackage{amsmath}
\usepackage{array}
\usepackage{longtable}
\usepackage{makecell}
\usepackage{booktabs}
\usepackage{caption}

\usepackage{parskip}
\usepackage{amsthm} 
\usepackage{xcolor}
\usepackage{pdflscape}
\usepackage{bm}

\usepackage{listings}

\usepackage[shortlabels]{enumitem}
\usepackage[margin=2cm]{geometry}

\newcommand{\dgeneral}{G}
\newcommand{\dtime}{T}
\newcommand{\dcost}{C}
\newcommand{\dquality}{Q}
\newcommand{\dflexibility}{F}

\newcommand{\ggroup}{G}
\newcommand{\gcase}{C}
\newcommand{\gact}{A}
\newcommand{\ginst}{I}

\newcommand{\universe}[2][]{U^{#1}_{#2}}
\newcommand{\attribute}[2][]{\#^{#1}_{#2}}

\newcommand{\atevent}{\mathit{event}}
\newcommand{\atcase}{\mathit{case}}
\newcommand{\atact}{\mathit{act}}
\newcommand{\atres}{\mathit{res}}
\newcommand{\athres}{\mathit{hres}}
\newcommand{\atnhres}{\mathit{nhres}}
\newcommand{\atrole}{\mathit{role}}
\newcommand{\attime}{\mathit{time}}
\newcommand{\attcost}{\mathit{tc}}
\newcommand{\atfcost}{\mathit{fc}}
\newcommand{\atvcost}{\mathit{vc}}
\newcommand{\atlcost}{\mathit{lc}}
\newcommand{\aticost}{\mathit{ic}}
\newcommand{\atmaincost}{\mathit{mainc}}
\newcommand{\atmdcost}{\mathit{mdc}}
\newcommand{\attranscost}{\mathit{transc}}
\newcommand{\atwarecost}{\mathit{warec}}
\newcommand{\atqual}{\mathit{qual}}
\newcommand{\atunit}{\mathit{unt}}
\newcommand{\atunsunit}{\mathit{uns}}
\newcommand{\atclient}{\mathit{cli}}
\newcommand{\atatt}{\mathit{att}}
\newcommand{\atval}{\mathit{val}}
\newcommand{\atmap}{\mathit{map}}
\newcommand{\attype}{\mathit{type}}
\newcommand{\atactinst}{\mathit{actinst}}
\newcommand{\ateventid}{\mathit{eid}}

\newcommand{\fnmin}{\mathit{min}}
\newcommand{\fnmax}{\mathit{max}}
\newcommand{\fnmatch}{\mathit{match}}
\newcommand{\fntoderintlog}{\mathit{\textbf{Der}}}
\newcommand{\fntototalderintlog}{\mathit{\textbf{TotDer}}}
\newcommand{\fntoexpintlog}{\mathit{\textbf{Expl}}}
\newcommand{\fnmatchone}{\mathit{match1}}
\newcommand{\fnmatchtwo}{\mathit{match2}}
\newcommand{\fnmatchfirst}{\mathit{matchfirst}}
\newcommand{\fnmatchlater}{\mathit{matchlater}}
\newcommand{\fnreassign}{\mathit{reassign}}
\newcommand{\fnseventbadness}{\mathit{sbadness}}
\newcommand{\fnceventbadness}{\mathit{cbadness}}
\newcommand{\fncompatible}{\mathit{cmp}}
\newcommand{\fnactinstcompatible}{\mathit{actinstcmp}}
\newcommand{\fnfresheventid}{\mathit{fresheventid}}
\newcommand{\fnfreshactinst}{\mathit{freshactinst}}
\newcommand{\fnevent}{\mathit{events}}
\newcommand{\fncase}{\mathit{case}}
\newcommand{\fnact}{\mathit{act}}
\newcommand{\fnres}{\mathit{res}}
\newcommand{\fnhres}{\mathit{hres}}
\newcommand{\fnrole}{\mathit{role}}
\newcommand{\fninst}{\mathit{inst}}
\newcommand{\fndom}{\mathit{dom}}
\newcommand{\fnleadtime}{\mathit{lt}}
\newcommand{\fnservicetime}{\mathit{st}}
\newcommand{\fnwaitingtime}{\mathit{wt}}
\newcommand{\fninstbetween}{\mathit{instbetween}}
\newcommand{\fntimewithin}{\mathit{timew}}
\newcommand{\fnstarttime}{\mathit{startt}}
\newcommand{\fnendtime}{\mathit{endt}}
\newcommand{\fnstartinstance}{\mathit{strin}}
\newcommand{\fnendinstance}{\mathit{endin}}
\newcommand{\fnistart}{\mathit{str}}
\newcommand{\fnicomplete}{\mathit{cpl}}
\newcommand{\fnistarttime}{\mathit{stime}}
\newcommand{\fnicompletetime}{\mathit{ctime}}
\newcommand{\fnprev}{\mathit{prev}}
\newcommand{\fnnext}{\mathit{next}}
\newcommand{\fnprevstr}{\mathit{prevstr}}
\newcommand{\fnconcstr}{\mathit{concstr}}
\newcommand{\fndres}{\mathit{dres}}
\newcommand{\fnfirst}{\mathit{fi}}
\newcommand{\fnfirstinsttc}{\mathit{fitc}}
\newcommand{\fnfirstinstlt}{\mathit{filt}}
\newcommand{\fndfr}{\mathit{dfrel}}
\newcommand{\fncount}{\mathit{count}}
\newcommand{\fnseq}{\mathit{seq}}
\newcommand{\fntrace}{\mathit{trace}}
\newcommand{\fnvariants}{\mathit{variants}}

\newcommand{\lsautl}{\mathit{Autl}}
\newcommand{\lsdesl}{\mathit{Desl}}
\newcommand{\lsdcl}{\mathit{DCl}}
\newcommand{\lsunwl}{\mathit{Unwl}}

\newcommand{\utif}{\mathit{if \ }}
\newcommand{\utand}{\mathit{\ and \ }}
\newcommand{\utwhere}{\mathit{\ where \ }}
\newcommand{\utwith}{\mathit{\ with \ }}
\newcommand{\utst}{\mid}
\newcommand{\utsuchthat}{\mathit{\ such~that \ }}
\newcommand{\utforall}{\mathit{for~all \ }}
\newcommand{\utforany}{\mathit{for~any \ }}
\newcommand{\utacta}{\mathit{a}}
\newcommand{\utactb}{\mathit{b}}
\newcommand{\uthres}{\mathit{hr}}
\newcommand{\utrole}{\mathit{rl}}
\newcommand{\utval}{\mathit{val}}
\newcommand{\utsubset}{\mathit{sub}}
\newcommand{\utstime}{\mathit{st}}
\newcommand{\utetime}{\mathit{et}}
\newcommand{\utexpected}{\mathit{e}}
\newcommand{\utwithin}{\mathit{^{w}}}
\newcommand{\utstarted}{\mathit{^{s}}}
\newcommand{\utcompleted}{\mathit{^{c}}}
\newcommand{\utstartedcompleted}{\mathit{^{sc}}}
\newcommand{\utsingle}{\mathit{^{sgl}}}
\newcommand{\utsummation}{\mathit{^{sum}}}
\newcommand{\utanyvariation}{\mathit{^{\times}}}
\newcommand{\utwithinnrm}{\mathit{w}}
\newcommand{\utstartednrm}{\mathit{s}}
\newcommand{\utcompletednrm}{\mathit{c}}
\newcommand{\utstartedcompletednrm}{\mathit{sc}}
\newcommand{\utsinglenrm}{\mathit{sgl}}
\newcommand{\utsummationnrm}{\mathit{sum}}
\newcommand{\utanyvariationnrm}{\mathit{\times}}

\newcommand{\atomiclog}{atomic event log}
\newcommand{\alog}{^{a}}
\newcommand{\derintlog}{derivable-interval event log}
\newcommand{\dilog}{^{di}}
\newcommand{\uniquederintlog}{uniquely-matched derivable-interval event log}
\newcommand{\totalderintlog}{totally-ordered derivable-interval event log}
\newcommand{\tdilog}{^{tdi}}
\newcommand{\expintlog}{explicit-interval event log}
\newcommand{\eilog}{^{ei}}
\newcommand{\transitlog}{transitional event log}
\newcommand{\tlog}{^{t}}

\newcommand{\inda}{\mathit{A}}
\newcommand{\indi}{\mathit{I}}
\newcommand{\indc}{\mathit{C}}
\newcommand{\inde}{\mathit{E}}
\newcommand{\indhr}{\mathit{HR}}
\newcommand{\indr}{\mathit{R}}
\newcommand{\indrole}{\mathit{Role}}

\newcommand{\indat}{\mathit{AT}}
\newcommand{\indauta}{\mathit{Auta}}
\newcommand{\indauti}{\mathit{Auti}}
\newcommand{\indautst}{\mathit{AutST}}
\newcommand{\indcltr}{\mathit{CLTR}}
\newcommand{\indcltvalc}{\mathit{CLTOC}}
\newcommand{\indcltvalp}{\mathit{CLTOP}}
\newcommand{\indcmdp}{\mathit{CMDP}}
\newcommand{\indh}{\mathit{H}}
\newcommand{\indit}{\mathit{IT}}
\newcommand{\indlt}{\mathit{LT}}
\newcommand{\indltcr}{\mathit{LTCR}}
\newcommand{\indltdl}{\mathit{LTDL}}
\newcommand{\indltde}{\mathit{LTDE}}
\newcommand{\indltfa}{\mathit{LTfA}}
\newcommand{\indltab}{\mathit{LTAB}}
\newcommand{\indltta}{\mathit{LTtA}}
\newcommand{\indsltr}{\mathit{SLTR}}
\newcommand{\indst}{\mathit{ST}}
\newcommand{\indstab}{\mathit{STAB}}
\newcommand{\indstu}{\mathit{STU}}
\newcommand{\indwt}{\mathit{WT}}
\newcommand{\indwtab}{\mathit{WTAB}}

\newcommand{\indautc}{\mathit{AutC}}
\newcommand{\inddac}{\mathit{DAC}}
\newcommand{\inddc}{\mathit{DC}}
\newcommand{\indfc}{\mathit{FC}}
\newcommand{\indhrcr}{\mathit{HRCR}}
\newcommand{\indic}{\mathit{IC}}
\newcommand{\indlctcr}{\mathit{LCTCR}}
\newcommand{\indlc}{\mathit{LC}}
\newcommand{\indmainc}{\mathit{MainC}}
\newcommand{\indmdc}{\mathit{MDC}}
\newcommand{\indoc}{\mathit{OC}}
\newcommand{\indrc}{\mathit{RC}}
\newcommand{\indrewc}{\mathit{RewC}}
\newcommand{\indrewp}{\mathit{RewP}}
\newcommand{\indtc}{\mathit{TC}}
\newcommand{\indtcltr}{\mathit{TCLTR}}
\newcommand{\indtcur}{\mathit{TCUR}}
\newcommand{\indtcstr}{\mathit{TCSTR}}
\newcommand{\indtransc}{\mathit{TransC}}
\newcommand{\indvc}{\mathit{VC}}
\newcommand{\indwarec}{\mathit{WareC}}

\newcommand{\indibyhr}{\mathit{IbyHR}}
\newcommand{\indibyrole}{\mathit{IbyRole}}
\newcommand{\indcaaftfc}{\mathit{CApastTFC}}
\newcommand{\indcabetfc}{\mathit{CAdueTFC}}
\newcommand{\indcaintfc}{\mathit{CAinTFC}}
\newcommand{\indcendac}{\mathit{CendAC}}
\newcommand{\indcstartac}{\mathit{CstartAC}}
\newcommand{\indcrewc}{\mathit{CRewC}}
\newcommand{\indcaaftfp}{\mathit{CApastTFP}}
\newcommand{\indcabetfp}{\mathit{CAdueTFP}}
\newcommand{\indcaintfp}{\mathit{CAinTFP}}
\newcommand{\indcendap}{\mathit{CendAP}}
\newcommand{\indcstartap}{\mathit{CstartAP}}
\newcommand{\indcrewp}{\mathit{CRewP}}
\newcommand{\indclitcr}{\mathit{CliTCR}}
\newcommand{\indcclir}{\mathit{CCliR}}
\newcommand{\indnauta}{\mathit{NAutA}}
\newcommand{\indnauti}{\mathit{NAutI}}
\newcommand{\indu}{\mathit{U}}
\newcommand{\indoq}{\mathit{OQ}}
\newcommand{\indrep}{\mathit{Rep}}
\newcommand{\indrewcv}{\mathit{RewCV}}
\newcommand{\indrewsub}{\mathit{RewSub}}
\newcommand{\indrewpv}{\mathit{RewPV}}
\newcommand{\indrt}{\mathit{RT}}
\newcommand{\indsuc}{\mathit{SUC}}
\newcommand{\indsup}{\mathit{SUP}}
\newcommand{\indtcclir}{\mathit{TCCliR}}
\newcommand{\induac}{\mathit{UAC}}
\newcommand{\induap}{\mathit{UAP}}
\newcommand{\induic}{\mathit{UIC}}
\newcommand{\induip}{\mathit{UIP}}

\newcommand{\indaroler}{\mathit{ARoleR}}
\newcommand{\indihrr}{\mathit{IHRR}}
\newcommand{\indcli}{\mathit{Cli}}
\newcommand{\indopta}{\mathit{OptA}}
\newcommand{\indopt}{\mathit{Opt}}
\newcommand{\inddfrar}{\mathit{DFRAR}}
\newcommand{\inddfr}{\mathit{DFR}}
\newcommand{\indrolevr}{\mathit{RoleVR}}
\newcommand{\indvcc}{\mathit{VCC}}
\newcommand{\indv}{\mathit{V}}

\newtheorem{definition}{Definition}
\newtheorem{proposition}{Proposition}
\newtheorem{theorem}{Theorem}
\newtheorem{lemma}{Lemma}

\setlength{\parskip}{\medskipamount}
\setlength{\parindent}{0pt}

\renewcommand{\emptyset}{\varnothing}

\usepackage{listings}
\lstset{
  language = Python,
  backgroundcolor = \color{lightgray!10},
  basicstyle = \ttfamily\small,
  keywordstyle = \color{blue}\bfseries,
  stringstyle = \color{red},
  commentstyle = \color{green!50!black},
  showstringspaces = false,
  numbers = left,
  numberstyle = \tiny\color{gray},
  frame = single,
  breaklines = true,
  captionpos = b,
  tabsize = 4,
  basicstyle = \footnotesize\ttfamily
}

\title{PPI Research - Definitions}
\author{Ignacio Velásquez}
\date{March 2026}

\begin{document}

\maketitle

\section{PPIs}

\footnotesize
\setlength\tabcolsep{2pt}
\setlength{\extrarowheight}{3pt}

\begin{longtable}{>{\hspace{0pt}}m{0.14\linewidth}>{\hspace{0pt}}c>{\hspace{0pt}}c>{\hspace{0pt}}m{0.1\linewidth}>{\hspace{0pt}}m{0.19\linewidth}>{\hspace{0pt}}m{0.473\linewidth}}
\toprule
\textbf{Indicator} & \textbf{D.} & \textbf{G.} & \textbf{F. Par.} & \textbf{Function Name} & \textbf{Formal Definition}\\
\midrule \endfirsthead
\toprule
\textbf{Indicator} & \textbf{D.} & \textbf{G.} & \textbf{F. Par.} & \textbf{Function Name} & \textbf{Formal Definition}\\
\midrule \endhead
\bottomrule \endfoot
\captionsetup{font=small}
\caption{General PPI definitions} \endlastfoot
Activity count (case granularity) & \dgeneral & \gcase & \(c \in C\) & \(\inda(c)\) & \(|\fnact(c)|\) \\
Activity count (group of cases granularity) & \dgeneral & \ggroup & \(g \subseteq C \colon\) \newline \(g \neq \emptyset\) & \(\inda(g)\) & \(|\cup_{c \in g} \fnact(c)|\) \\
Expected Activity count & \dgeneral & \ggroup & \(g \subseteq C \colon\) \newline \(g \neq \emptyset\) & \(\utexpected\inda(g)\) & \(\frac{\sum _{c \in g} \inda(c)}{\indc(g)}\) \\
activity Instance count (activity granularity) & \dgeneral & \gact & \(a \in A\) & \(\indi(a)\) & \(|\fninst(a)|\) \\
activity Instance count (case granularity) & \dgeneral & \gcase & \(c \in C\) & \(\indi(c)\) & \(|\fninst(c)|\) \\
activity Instance count (group of cases granularity) & \dgeneral & \ggroup & \(g \subseteq C \colon\) \newline \(g \neq \emptyset\) & \(\indi(g)\) & \(\sum _{c \in g} \indi(c)\) \\
Expected activity Instance count & \dgeneral & \ggroup & \(g \subseteq C \colon\) \newline \(g \neq \emptyset\) & \(\utexpected\indi(g)\) & \(\frac{\indi(g)}{\indc(g)}\) \\
Case count (group of cases granularity) & \dgeneral & \ggroup & \(g \subseteq C \colon\) \newline \(g \neq \emptyset\) & \(\indc(g)\) & \(|g|\) \\
Human Resource count (activity granularity) & \dgeneral & \gact & \(a \in A\) & \(\indhr(a)\) & \(|\fnhres(a)|\) \\
Human Resource count (case granularity) & \dgeneral & \gcase & \(c \in C\) & \(\indhr(c)\) & \(|\fnhres(c)|\) \\
Human Resource count (group of cases granularity) & \dgeneral & \ggroup & \(g \subseteq C \colon\) \newline \(g \neq \emptyset\) & \(\indhr(g)\) & \(|\cup_{c \in g} \fnhres(c)|\) \\
Expected Human Resource count & \dgeneral & \ggroup & \(g \subseteq C \colon\) \newline \(g \neq \emptyset\) & \(\utexpected\indhr(g)\) & \(\frac{\sum _{c \in g} \indhr(c)}{\indc(g)}\) \\
Resource count (activity granularity) & \dgeneral & \gact & \(a \in A\) & \(\indr(a)\) & \(|\fnres(a)|\) \\
Resource count (case granularity) & \dgeneral & \gcase & \(c \in C\) & \(\indr(c)\) & \(|\fnres(c)|\) \\
Resource count (group of cases granularity) & \dgeneral & \ggroup & \(g \subseteq C \colon\) \newline \(g \neq \emptyset\) & \(\indr(g)\) & \(|\cup_{c \in g} \fnres(c)|\) \\
Expected Resource count & \dgeneral & \ggroup & \(g \subseteq C \colon\) \newline \(g \neq \emptyset\) & \(\utexpected\indr(g)\) & \(\frac{\sum _{c \in g} \indr(c)}{\indc(g)}\) \\
Role count (case granularity) & \dgeneral & \gcase & \(c \in C\) & \(\indrole(c)\) & \(|\fnrole(c)|\) \\
Role count (group of cases granularity) & \dgeneral & \ggroup & \(g \subseteq C \colon\) \newline \(g \neq \emptyset\) & \(\indrole(g)\) & \(|\cup_{c \in g} \fnrole(c)|\) \\
Expected Role count & \dgeneral & \ggroup & \(g \subseteq C \colon\) \newline \(g \neq \emptyset\) & \(\utexpected\indrole(g)\) & \(\frac{\sum _{c \in g} \indrole(c)}{\indc(g)}\) \\
\bottomrule
\label{PPIGeneralDefinitions}
\end{longtable}


\begin{longtable}{>{\hspace{0pt}}m{0.14\linewidth}>{\hspace{0pt}}c>{\hspace{0pt}}c>{\hspace{0pt}}m{0.1\linewidth}>{\hspace{0pt}}m{0.19\linewidth}>{\hspace{0pt}}m{0.473\linewidth}}
\toprule
\textbf{Indicator} & \textbf{D.} & \textbf{G.} & \textbf{F. Par.} & \textbf{Function Name} & \textbf{Formal Definition}\\
\midrule \endfirsthead
\toprule
\textbf{Indicator} & \textbf{D.} & \textbf{G.} & \textbf{F. Par.} & \textbf{Function Name} & \textbf{Formal Definition}\\
\midrule \endhead
\bottomrule \endfoot
\captionsetup{font=small}
\caption{Time dimension PPI definitions} \endlastfoot
Active Time (case granularity) & \dtime & \gcase & \(c \in C\) & \(\indat(c)\) & \(\indlt(c) - \indit(c)\) \\
Expected Active Time & \dtime & \ggroup & \(g \subseteq C \colon\) \newline \(g \neq \emptyset\) & \(\utexpected\indat(g)\) & \(\frac{\sum _{c \in g} \indat(c)}{\indc(g)}\) \\
Activity count (case granularity) & \dtime & \gcase & \(c \in C\) & \(\inda(c)\) & \(|\fnact(c)|\) \\
Activity count (group of cases granularity) & \dtime & \ggroup & \(g \subseteq C \colon\) \newline \(g \neq \emptyset\) & \(\inda(g)\) & \(|\cup_{c \in g} \fnact(c)|\) \\
Expected Activity count & \dtime & \ggroup & \(g \subseteq C \colon\) \newline \(g \neq \emptyset\) & \(\utexpected\inda(g)\) & \(\frac{\sum _{c \in g} \inda(c)}{\indc(g)}\) \\
activity Instance count (case granularity) & \dtime & \gcase & \(c \in C\) & \(\indi(c)\) & \(|\fninst(c)|\) \\
activity Instance count (group of cases granularity) & \dtime & \ggroup & \(g \subseteq C \colon\) \newline \(g \neq \emptyset\) & \(\indi(g)\) & \(\sum _{c \in g} \indi(c)\) \\
Expected activity Instance count & \dtime & \ggroup & \(g \subseteq C \colon\) \newline \(g \neq \emptyset\) & \(\utexpected\indi(g)\) & \(\frac{\indi(g)}{\indc(g)}\) \\
Automated Activity count (case granularity) & \dtime & \gcase & \(c \in C\) & \(\indauta(c)\) & \(|\lsautl \cap \fnact(c)|\) \\
Automated Activity count (group of cases granularity) & \dtime & \ggroup & \(g \subseteq C \colon\) \newline \(g \neq \emptyset\) & \(\indauta(g)\) & \(|\lsautl \cap (\cup_{c \in g} \fnact(c))|\) \\
Expected Automated Activity count & \dtime & \ggroup & \(g \subseteq C \colon\) \newline \(g \neq \emptyset\) & \(\utexpected\indauta(g)\) & \(\frac{\sum _{c \in g} \indauta(c)}{\indc(g)}\) \\
Automated activity Instance count (case granularity) & \dtime & \gcase & \(c \in C\) & \(\indauti(c)\) & \(|\{i \in \fninst(c) \utst \fnact(i) \in \lsautl\}|\) \\
Automated activity Instance count (group of cases granularity) & \dtime & \ggroup & \(g \subseteq C \colon\) \newline \(g \neq \emptyset\) & \(\indauti(g)\) & \(\sum _{c \in g} \indauti(c)\) \\
Expected Automated activity Instance count & \dtime & \ggroup & \(g \subseteq C \colon\) \newline \(g \neq \emptyset\) & \(\utexpected\indauti(g)\) & \(\frac{\indauti(g)}{\indc(g)}\) \\
Automated Activity Service Time (case granularity) & \dtime & \gcase & \(c \in C\) & \(\indautst(c)\) & \(\sum _{i \in \fninst(c), \fnact(i) \in \lsautl} \indst(i)\) \\
Automated Activity Service Time (group of cases granularity) & \dtime & \ggroup & \(g \subseteq C \colon\) \newline \(g \neq \emptyset\) & \(\indautst(g)\) & \(\sum _{c \in g} \indautst(c)\) \\
Expected Automated Activity Service Time & \dtime & \ggroup & \(g \subseteq C \colon\) \newline \(g \neq \emptyset\) & \(\utexpected\indautst(g)\) & \(\frac{\indautst(g)}{\indc(g)}\) \\
Case count and Lead Time Ratio (group of cases granularity) & \dtime & \ggroup & \(g \subseteq C \colon\) \newline \(g \neq \emptyset\) & \(\indcltr(g)\) & \(\frac{\indc(g)}{\indlt(g)}\) \\
Case Count where Lead Time Over value (group of cases granularity) & \dtime & \ggroup & \(g \subseteq C \colon\) \newline \(g \neq \emptyset\) \newline \(\utval \in \mathbb{R}\) & \(\indcltvalc(g, \utval)\) & \(|\{c \in g \utst \indlt(c) > \utval\}|\) \\
Case Percentage where Lead Time Over value (group of cases granularity) & \dtime & \ggroup & \(g \subseteq C \colon\) \newline \(g \neq \emptyset\) \newline \(\utval \in \mathbb{R}\) & \(\indcltvalp(g, \utval)\) & \(\frac{\indcltvalc(g, \utval)}{\indc(g)}\) \\
Case Percentage with Missed Deadline (group of cases granularity) & \dtime & \ggroup & \(g \subseteq C \colon\) \newline \(g \neq \emptyset\) \newline \(\utval \in \universe{\attime}\) & \(\indcmdp(g, \utval)\) & \(\frac{|\{c \in g \utst \fnendtime(c) > \utval\}|}{\indc(g)}\) \\
Handover count (case granularity) & \dtime & \gcase & \(c \in C\) & \(\indh(c)\) & \(\sum _{i \in \fninst(c)} \fndres(i)\) \\
Expected Handover count & \dtime & \ggroup & \(g \subseteq C \colon\) \newline \(g \neq \emptyset\) & \(\utexpected\indh(g)\) & \(\frac{\sum _{c \in g} \indh(c)}{\indc(g)}\) \\
Idle Time (case granularity) & \dtime & \gcase & \(c \in C\) & \(\indit(c)\) & \(\sum _{i \in \fninst(c)} \begin{cases}
\frac{\indwt(i)}{|\fnconcstr(i)|} & \utif \fnprevstr(i) = \emptyset \\
0 & \utif \fnprevstr(i) \neq \emptyset
\end{cases}\) \\
Expected Idle Time & \dtime & \ggroup & \(g \subseteq C \colon\) \newline \(g \neq \emptyset\) & \(\utexpected\indit(g)\) & \(\frac{\sum _{c \in g} \indit(c)}{\indc(g)}\) \\
Lead Time (activity instance granularity) & \dtime & \ginst & \(i \in I\) & \(\indlt(i)\) & \(\indst(i) + \indwt(i)\) \\
Lead Time (activity granularity) & \dtime & \gact & \(a \in A\) & \(\indlt(a)\) & \(\sum _{i \in \fninst(a)} \indlt(i)\) \\
Lead Time (case granularity) & \dtime & \gcase & \(c \in C\) & \(\indlt(c)\) & \(\fnendtime(c) - \fnstarttime(c)\) \\
Lead Time (group of cases granularity) & \dtime & \ggroup & \(g \subseteq C \colon\) \newline \(g \neq \emptyset\)& \(\indlt(g)\) & \(^{\fnmax}_{c \in g} \fnendtime(c) - ^{\fnmin}_{c \in g} \fnstarttime(c)\) \\
Expected Lead Time & \dtime & \ggroup & \(g \subseteq C \colon\) \newline \(g \neq \emptyset\) & \(\utexpected\indlt(g)\) & \(\frac{\sum _{c\in g}\indlt(c)}{\indc(g)}\) \\
Lead Time and Case count Ratio (group of cases granularity) & \dtime & \ggroup & \(g \subseteq C \colon\) \newline \(g \neq \emptyset\) & \(\indltcr(g)\) & \(\frac{\indlt(g)}{\indc(g)}\) \\
Lead Time Deviation from time Limit (case granularity) & \dtime & \gcase & \(c \in C\) \newline \(\utval \in \mathbb{R}\) & \(\indltdl(c, \utval)\) & \(\utval - \indlt(c)\) \\
Expected Lead Time Deviation from time Limit & \dtime & \ggroup & \(g \subseteq C \colon\) \newline \(g \neq \emptyset\) \newline \(\utval \in \mathbb{R}\) & \(\utexpected\indltdl(g, \utval)\) & \(\frac{\sum _{c \in g} \indltdl(c, \utval)}{\indc(g)}\) \\
Lead Time Deviation from Expectation (case granularity) & \dtime & \gcase & \(c \in C\) \newline \(\utval \in \mathbb{R}\) & \(\indltde(c, \utval)\) & \(|\utval - \indlt(c)|\) \\
Expected Lead Time Deviation from Expectation & \dtime & \ggroup & \(g \subseteq C \colon\) \newline \(g \neq \emptyset\) \newline \(\utval \in \mathbb{R}\) & \(\utexpected\indltde(g, \utval)\) & \(\frac{\sum _{c \in g} \indltde(c, \utval)}{\indc(g)}\) \\
Lead Time from activity A (case granularity) & \dtime & \gcase & \(c \in C\) \newline \(\utacta \in A\) & \(\indltfa(c, \utacta)\) & \(\begin{cases}
\fnleadtime(x, y) & \utforany x \in \fnfirst\utstarted(c, \utacta) \utand y \in \fnendinstance(c) \\ \bot & \utif \fnfirst\utstarted(c, \utacta) = \emptyset
\end{cases}\) \\
Expected Lead Time from activity A & \dtime & \ggroup & \(g \subseteq C \colon\) \newline \(g \neq \emptyset\) \newline \(\utacta \in A\) & \(\utexpected\indltfa(g, \utacta)\) & \(\frac{\sum _{c \in g} \begin{cases}
\indltfa(c, \utacta) & \utif \fnfirst\utstarted(c, \utacta) \neq \emptyset \\ 0 & \utif \fnfirst\utstarted(c, \utacta) = \emptyset
\end{cases}}{|\{c \in g \utst \fnfirst\utstarted(c, \utacta) \neq \emptyset\}|}\) \\
Lead Time from activity A to activity B (case granularity) & \dtime & \gcase & \(c \in C\) \newline \(\utacta, \utactb \in A \colon\) \newline \(\utacta \neq \utactb\) & \(\indltab(c, \utacta, \utactb)\) & \(\begin{cases}
\fnleadtime(x,y) & \utforany x \in \fnfirst\utstarted(c, \utacta) \utand y \in \fnfirst\utstartedcompleted(c, \utacta, \utactb) \\
\bot & \utif \fnfirst\utstartedcompleted(c, \utacta, \utactb) = \emptyset
\end{cases}\) \\
Expected Lead Time from activity A to activity B & \dtime & \ggroup & \(g \subseteq C \colon\) \newline \(g \neq \emptyset\) \newline \(\utacta, \utactb \in A \colon\) \newline \(\utacta \neq \utactb\) & \(\utexpected\indltab(g, \utacta, \utactb)\) & \(\frac{\sum _{c \in g} \begin{cases}
\indltab(c, \utacta, \utactb) & \utif \fnfirst\utstartedcompleted(c, \utacta, \utactb) \neq \emptyset \\ 0 & \utif \fnfirst\utstartedcompleted(c, \utacta, \utactb) = \emptyset
\end{cases}}{|\{c \in g \utst \fnfirst\utstartedcompleted(c, \utacta, \utactb) \neq \emptyset)\}|}\) \\
Lead Time to activity A (case granularity) & \dtime & \gcase & \(c \in C\) \newline \(\utacta \in A\) & \(\indltta(c, \utacta)\) & \(\begin{cases}
\fnleadtime(x, y) & \utforany x \in \fnstartinstance(c) \utand y \in \fnfirst\utcompleted(c, \utacta) \\ \bot & \utif \fnfirst\utcompleted(c, \utacta) = \emptyset
\end{cases}\) \\
Expected Lead Time to activity A & \dtime & \ggroup & \(g \subseteq C \colon\) \newline \(g \neq \emptyset\) \newline \(\utacta \in A\) & \(\utexpected\indltta(g, \utacta)\) & \(\frac{\sum _{c \in g} \begin{cases}
\indltta(c, \utacta) & \utif \fnfirst\utcompleted(c, \utacta) \neq \emptyset \\ 0 & \utif \fnfirst\utcompleted(c, \utacta) = \emptyset
\end{cases}}{|\{c \in g \utst \utif \fnfirst\utcompleted(c, \utacta) \neq \emptyset\}|}\) \\
Service and Lead Time Ratio (activity instance granularity) & \dtime & \ginst & \(i \in I\) & \(\indsltr(i)\) & \(\frac{\indst(i)}{\indlt(i)}\) \\
Service and Lead Time Ratio (activity granularity) & \dtime & \gact & \(a \in A\) & \(\indsltr(a)\) & \(\frac{\indst(a)}{\indlt(a)}\) \\
Service and Lead Time Ratio (case granularity) & \dtime & \gcase & \(c \in C\) & \(\indsltr(c)\) & \(\frac{\indst(c)}{\indlt(c)}\) \\
Service and Lead Time Ratio (group of cases granularity) & \dtime & \ggroup & \(g \subseteq C \colon\) \newline \(g \neq \emptyset\) & \(\indsltr(g)\) & \(\frac{\indst(g)}{\indlt(g)}\) \\
Expected Service and Lead Time Ratio & \dtime & \ggroup & \(g \subseteq C \colon\) \newline \(g \neq \emptyset\) & \(\utexpected\indsltr(g)\) & \(\frac{\indst(g)}{\sum _{c \in g} \indlt(c)}\) \\
Service Time (activity instance granularity) & \dtime & \ginst & \(i \in I\) & \(\indst(i)\) & \(\fnicompletetime(i) - \fnistarttime(i)\) \\
Service Time (activity granularity) & \dtime & \gact & \(a \in A\) & \(\indst(a)\) & \(\sum _{i \in \fninst(a)} \indst(i)\) \\
Service Time (case granularity) & \dtime & \gcase & \(c \in C\) & \(\indst(c)\) & \(\sum _{i \in \fninst(c)} \indst(i)\) \\
Service Time (group of cases granularity) & \dtime & \ggroup & \(g \subseteq C \colon\) \newline \(g \neq \emptyset\) & \(\indst(g)\) & \(\sum _{c \in g} \indst(c)\) \\
Expected Service Time & \dtime & \ggroup & \(g \subseteq C \colon\) \newline \(g \neq \emptyset\) & \(\utexpected\indst(g)\) & \(\frac{\indst(g)}{\indc(g)}\) \\
Service Time from activity A to activity B (case granularity) & \dtime & \gcase & \(c \in C\) \newline \(\utacta, \utactb \in A \colon\) \newline \(\utacta \neq \utactb\) & \(\indstab\utanyvariation(c, \utacta, \utactb)\)
\newline \newline \(\utanyvariationnrm \in \{\utstartednrm, \utcompletednrm, \utstartedcompletednrm, \utwithinnrm\}\) & \(\begin{cases} \fnservicetime\utanyvariation(x, y) & \utforany x \in \fnfirst\utstarted(c, \utacta) \utand y \in \fnfirst\utstartedcompleted(c, \utacta, \utactb) \\
\bot & \utif \fnfirst\utstartedcompleted(c, \utacta, \utactb) = \emptyset \end{cases}\)
\newline \newline The $\utanyvariationnrm$ indicates that the PPI can take multiple forms, in this case, $\utanyvariationnrm \in \{\utstartednrm, \utcompletednrm, \utstartedcompletednrm, \utwithinnrm\}$. If $\utanyvariationnrm = \utstartednrm$, then the function considers activity instances that were started within the start and end activity instances, i.e., it uses $\fnservicetime\utstarted$; if $\utanyvariationnrm = \utcompletednrm$, then the function considers activity instances that were completed within the start and end activity instances, i.e., it uses $\fnservicetime\utcompleted$; if $\utanyvariationnrm = \utstartedcompletednrm$, then the function considers activity instances that were either started or completed within the start and end activity instances, i.e., it uses $\fnservicetime\utstartedcompleted$; if $\utanyvariationnrm = \utwithin$, then the function considers all activity instances that were active within the start and end activity instances, i.e., it uses $\fnservicetime\utwithin$. \\
Service Time from activity A to activity B (group of cases granularity) & \dtime & \ggroup & \(g \subseteq C \colon\) \newline \(g \neq \emptyset\) \newline \(\utacta, \utactb \in A \colon\) \newline \(\utacta \neq \utactb\) & \(\indstab\utanyvariation(g, \utacta, \utactb)\)
\newline \newline \(\utanyvariationnrm \in \{\utstartednrm, \utcompletednrm, \utstartedcompletednrm, \utwithinnrm\}\) & \(\sum _{c \in g} \begin{cases}
\indstab\utanyvariation(c, \utacta, \utactb) & \utif \fnfirst\utstartedcompleted(c, \utacta, \utactb) \neq \emptyset \\ 0 & \utif \fnfirst\utstartedcompleted(c, \utacta, \utactb) = \emptyset
\end{cases}\)
\newline \newline The $\utanyvariationnrm$ indicates that the PPI can take multiple forms, in this case, $\utanyvariationnrm \in \{\utstartednrm, \utcompletednrm, \utstartedcompletednrm, \utwithinnrm\}$. If $\utanyvariationnrm = \utstartednrm$, then the function considers activity instances that were started within the start and end activity instances, i.e., it uses $\indstab\utstarted$; if $\utanyvariationnrm = \utcompletednrm$, then the function considers activity instances that were completed within the start and end activity instances, i.e., it uses $\indstab\utcompleted$; if $\utanyvariationnrm = \utstartedcompletednrm$, then the function considers activity instances that were either started or completed within the start and end activity instances, i.e., it uses $\indstab\utstartedcompleted$; if $\utanyvariationnrm = \utwithin$, then the function considers all activity instances that were active within the start and end activity instances, i.e., it uses $\indstab\utwithin$. \\
Expected Service Time from activity A to activity B & \dtime & \ggroup & \(g \subseteq C \colon\) \newline \(g \neq \emptyset\) \newline \(\utacta, \utactb \in A \colon\) \newline \(\utacta \neq \utactb\) & \(\utexpected\indstab\utanyvariation(g, \utacta, \utactb)\)
\newline \newline \(\utanyvariationnrm \in \{\utstartednrm, \utcompletednrm, \utstartedcompletednrm, \utwithinnrm\}\) & \(\frac{\indstab\utanyvariation(g, \utacta, \utactb)}{|\{c \in g \utst \fnfirst\utstartedcompleted(c, \utacta, \utactb) \neq \emptyset\}|}\)
\newline \newline The $\utanyvariationnrm$ indicates that the PPI can take multiple forms, in this case, $\utanyvariationnrm \in \{\utstartednrm, \utcompletednrm, \utstartedcompletednrm, \utwithinnrm\}$. If $\utanyvariationnrm = \utstartednrm$, then the function considers activity instances that were started within the start and end activity instances, i.e., it uses $\indstab\utstarted$; if $\utanyvariationnrm = \utcompletednrm$, then the function considers activity instances that were completed within the start and end activity instances, i.e., it uses $\indstab\utcompleted$; if $\utanyvariationnrm = \utstartedcompletednrm$, then the function considers activity instances that were either started or completed within the start and end activity instances, i.e., it uses $\indstab\utstartedcompleted$; if $\utanyvariationnrm = \utwithin$, then the function considers all activity instances that were active within the start and end activity instances, i.e., it uses $\indstab\utwithin$. \\
Waiting Time (activity instance granularity) & \dtime & \ginst & \(i \in I\) & \(\indwt(i)\) & \(\begin{cases}
\fnistarttime(i) - \fnicompletetime(x) & \utforany x \in \fnprev(i) \\
0 & \utif \fnprev(i) = \emptyset
\end{cases}\) \\
Waiting Time (activity granularity) & \dtime & \gact & \(a \in A\) & \(\indwt(a)\) & \(\sum _{i \in \fninst(a)} \indwt(i)\) \\
Waiting Time (case granularity) & \dtime & \gcase & \(c \in C\) & \(\indwt(c)\) & \(\sum _{i \in \fninst(c)} \indwt(i)\) \\
Expected Waiting Time & \dtime & \ggroup & \(g \subseteq C \colon\) \newline \(g \neq \emptyset\) & \(\utexpected\indwt(g)\) & \(\frac{\sum _{c \in g} \indwt(c)}{\indc(g)}\) \\
Waiting Time from activity A to activity B (case granularity) & \dtime & \gcase & \(c \in C\) \newline \(\utacta, \utactb \in A \colon\) \newline \(\utacta \neq \utactb\) & \(\indwtab\utanyvariation(c, \utacta, \utactb)\)
\newline \newline \(\utanyvariationnrm \in \{\utstartednrm, \utcompletednrm, \utstartedcompletednrm, \utwithinnrm\}\) & \(\begin{cases} \fnwaitingtime\utanyvariation(x, y) & \utforany x \in \fnfirst\utstarted(c, \utacta) \utand y \in \fnfirst\utstartedcompleted(c, \utacta, \utactb) \\
\bot & \utif \fnfirst\utstartedcompleted(c, \utacta, \utactb) = \emptyset \end{cases}\)
\newline \newline The $\utanyvariationnrm$ indicates that the PPI can take multiple forms, in this case, $\utanyvariationnrm \in \{\utstartednrm, \utcompletednrm, \utstartedcompletednrm, \utwithinnrm\}$. If $\utanyvariationnrm = \utstartednrm$, then the function considers activity instances that were started within the start and end activity instances, i.e., it uses $\fnwaitingtime\utstarted$; if $\utanyvariationnrm = \utcompletednrm$, then the function considers activity instances that were completed within the start and end activity instances, i.e., it uses $\fnwaitingtime\utcompleted$; if $\utanyvariationnrm = \utstartedcompletednrm$, then the function considers activity instances that were either started or completed within the start and end activity instances, i.e., it uses $\fnwaitingtime\utstartedcompleted$; if $\utanyvariationnrm = \utwithin$, then the function considers all activity instances that were active within the start and end activity instances, i.e., it uses $\fnwaitingtime\utwithin$. \\
Expected Waiting Time from activity A to activity B & \dtime & \ggroup & \(g \subseteq C \colon\) \newline \(g \neq \emptyset\) \newline \(\utacta, \utactb \in A \colon\) \newline \(\utacta \neq \utactb\) & \(\utexpected\indwtab\utanyvariation(g, \utacta, \utactb)\)
\newline \newline \(\utanyvariationnrm \in \{\utstartednrm, \utcompletednrm, \utstartedcompletednrm, \utwithinnrm\}\) & \(\frac{\sum _{c \in g} \begin{cases}
\indwtab\utanyvariation(c, \utacta, \utactb) & \utif \fnfirst\utstartedcompleted(c, \utacta, \utactb) \neq \emptyset \\ 0 & \utif \fnfirst\utstartedcompleted(c, \utacta, \utactb) = \emptyset
\end{cases}}{|\{c \in g \utst \fnfirst\utstartedcompleted(c, \utacta, \utactb) \neq \emptyset\}|}\)
\newline \newline The $\utanyvariationnrm$ indicates that the PPI can take multiple forms, in this case, $\utanyvariationnrm \in \{\utstartednrm, \utcompletednrm, \utstartedcompletednrm, \utwithinnrm\}$. If $\utanyvariationnrm = \utstartednrm$, then the function considers activity instances that were started within the start and end activity instances, i.e., it uses $\indwtab\utstarted$; if $\utanyvariationnrm = \utcompletednrm$, then the function considers activity instances that were completed within the start and end activity instances, i.e., it uses $\indwtab\utcompleted$; if $\utanyvariationnrm = \utstartedcompletednrm$, then the function considers activity instances that were either started or completed within the start and end activity instances, i.e., it uses $\indwtab\utstartedcompleted$; if $\utanyvariationnrm = \utwithin$, then the function considers all activity instances that were active within the start and end activity instances, i.e., it uses $\indwtab\utwithin$. \\
\bottomrule
\label{PPITimeDefinitions}
\end{longtable}


\begin{longtable}{>{\hspace{0pt}}m{0.14\linewidth}>{\hspace{0pt}}c>{\hspace{0pt}}c>{\hspace{0pt}}m{0.1\linewidth}>{\hspace{0pt}}m{0.19\linewidth}>{\hspace{0pt}}m{0.473\linewidth}}
\toprule
\textbf{Indicator} & \textbf{D.} & \textbf{G.} & \textbf{F. Par.} & \textbf{Function Name} & \textbf{Formal Definition}\\
\midrule \endfirsthead
\toprule
\textbf{Indicator} & \textbf{D.} & \textbf{G.} & \textbf{F. Par.} & \textbf{Function Name} & \textbf{Formal Definition}\\
\midrule \endhead
\bottomrule \endfoot
\captionsetup{font=small}
\caption{Cost dimension PPI definitions} \endlastfoot
Automated Activity Cost (case granularity) & \dcost & \gcase & \(c \in C\) & \(\indautc\utanyvariation(c)\)
\newline \newline \(\utanyvariationnrm \in \{\utsinglenrm, \utsummationnrm\}\) & \(\sum _{i \in \fninst(c), \fnact(i) \in \lsautl} \indtc\utanyvariation(i)\)
\newline \newline The $\utanyvariationnrm$ indicates that the PPI can take multiple forms, in this case, $\utanyvariationnrm \in \{\utsinglenrm, \utsummationnrm\}$. If $\utanyvariationnrm = \utsinglenrm$, then the function considers single events of activity instances for cost calculations, i.e., it uses $\indtc\utsingle$; if $\utanyvariationnrm = \utsummationnrm$, then the function considers the sum of all events of activity instances for cost calculations, i.e., it uses $\indtc\utsummation$. \\
Automated Activity Cost (group of cases granularity) & \dcost & \ggroup & \(g \subseteq C \colon\) \newline \(g \neq \emptyset\) & \(\indautc\utanyvariation(g)\)
\newline \newline \(\utanyvariationnrm \in \{\utsinglenrm, \utsummationnrm\}\) & \(\sum _{c \in g} \indautc\utanyvariation(c)\)
\newline \newline The $\utanyvariationnrm$ indicates that the PPI can take multiple forms, in this case, $\utanyvariationnrm \in \{\utsinglenrm, \utsummationnrm\}$. If $\utanyvariationnrm = \utsinglenrm$, then the function considers single events of activity instances for cost calculations, i.e., it uses $\indautc\utsingle$; if $\utanyvariationnrm = \utsummationnrm$, then the function considers the sum of all events of activity instances for cost calculations, i.e., it uses $\indautc\utsummation$. \\
Expected Automated Activity Cost & \dcost & \ggroup & \(g \subseteq C \colon\) \newline \(g \neq \emptyset\) & \(\utexpected\indautc\utanyvariation(g)\) 
\newline \newline \(\utanyvariationnrm \in \{\utsinglenrm, \utsummationnrm\}\) & \(\frac{\indautc\utanyvariation(g)}{\indc(g)}\)
\newline \newline The $\utanyvariationnrm$ indicates that the PPI can take multiple forms, in this case, $\utanyvariationnrm \in \{\utsinglenrm, \utsummationnrm\}$. If $\utanyvariationnrm = \utsinglenrm$, then the function considers single events of activity instances for cost calculations, i.e., it uses $\indautc\utsingle$; if $\utanyvariationnrm = \utsummationnrm$, then the function considers the sum of all events of activity instances for cost calculations, i.e., it uses $\indautc\utsummation$. \\
Desired Activity count (case granularity) & \dcost & \gcase & \(c \in C\) & \(\inddac(c)\) & \(| \lsdesl \cap \fnact(c)|\) \\
Desired Activity count (group of cases granularity) & \dcost & \ggroup & \(g \subseteq C \colon\) \newline \(g \neq \emptyset\) & \(\inddac(g)\) & \(|\lsdesl \cap (\cup_{c \in g} \fnact(c))|\) \\
Expected Desired Activity count & \dcost & \ggroup & \(g \subseteq C \colon\) \newline \(g \neq \emptyset\) & \(\utexpected\inddac(g)\) & \(\frac{\sum _{c \in g} \inddac(c)}{\indc(g)}\) \\
Direct Cost (case granularity) & \dcost & \gcase & \(c \in C\) & \(\inddc\utanyvariation(c)\) 
\newline \newline \(\utanyvariationnrm \in \{\utsinglenrm, \utsummationnrm\}\) & \(\sum _{i \in \fninst(c), \fnact(i) \in \lsdcl} \indtc\utanyvariation(i)\)
\newline \newline The $\utanyvariationnrm$ indicates that the PPI can take multiple forms, in this case, $\utanyvariationnrm \in \{\utsinglenrm, \utsummationnrm\}$. If $\utanyvariationnrm = \utsinglenrm$, then the function considers single events of activity instances for cost calculations, i.e., it uses $\indtc\utsingle$; if $\utanyvariationnrm = \utsummationnrm$, then the function considers the sum of all events of activity instances for cost calculations, i.e., it uses $\indtc\utsummation$. \\
Direct Cost (group of cases granularity) & \dcost & \ggroup & \(g \subseteq C \colon\) \newline \(g \neq \emptyset\) & \(\inddc\utanyvariation(g)\) 
\newline \newline \(\utanyvariationnrm \in \{\utsinglenrm, \utsummationnrm\}\) & \(\sum _{c \in g} \inddc\utanyvariation(c)\)
\newline \newline The $\utanyvariationnrm$ indicates that the PPI can take multiple forms, in this case, $\utanyvariationnrm \in \{\utsinglenrm, \utsummationnrm\}$. If $\utanyvariationnrm = \utsinglenrm$, then the function considers single events of activity instances for cost calculations, i.e., it uses $\inddc\utsingle$; if $\utanyvariationnrm = \utsummationnrm$, then the function considers the sum of all events of activity instances for cost calculations, i.e., it uses $\inddc\utsummation$. \\
Expected Direct Cost & \dcost & \ggroup & \(g \subseteq C \colon\) \newline \(g \neq \emptyset\) & \(\utexpected\inddc\utanyvariation(g)\) 
\newline \newline \(\utanyvariationnrm \in \{\utsinglenrm, \utsummationnrm\}\) & \(\frac{\inddc\utanyvariation(g)}{\indc(g)}\)
\newline \newline The $\utanyvariationnrm$ indicates that the PPI can take multiple forms, in this case, $\utanyvariationnrm \in \{\utsinglenrm, \utsummationnrm\}$. If $\utanyvariationnrm = \utsinglenrm$, then the function considers single events of activity instances for cost calculations, i.e., it uses $\inddc\utsingle$; if $\utanyvariationnrm = \utsummationnrm$, then the function considers the sum of all events of activity instances for cost calculations, i.e., it uses $\inddc\utsummation$. \\
Fixed Cost considering single events of activity instances (activity instance granularity) & \dcost & \ginst & \(i \in I\) & \(\indfc\utsingle(i)\) & \(\begin{cases}
\attribute{\atfcost}(\fnicomplete(i)) & \utif \attribute{\atfcost}(\fnicomplete(i)) \neq \bot \\
\attribute{\atfcost}(\fnistart(i)) & \utif \attribute{\atfcost}(\fnicomplete(i)) = \bot \land \attribute{\atfcost}(\fnistart(i)) \neq \bot \\
\bot & \utif \attribute{\atfcost}(\fnicomplete(i)) = \bot \land \attribute{\atfcost}(\fnistart(i)) = \bot \\
\end{cases}\) \\
Fixed Cost considering the sum of all events of activity instances (activity instance granularity) & \dcost & \ginst & \(i \in I\) & \(\indfc\utsummation(i)\) & \(\begin{cases}
\attribute{\atfcost}(\fnistart(i)) + & \\ \attribute{\atfcost}(\fnicomplete(i)) & \utif \attribute{\atfcost}(\fnicomplete(i)) \neq \bot \land \attribute{\atfcost}(\fnistart(i)) \neq \bot \\
\bot & \utif \attribute{\atfcost}(\fnicomplete(i)) = \bot \lor \attribute{\atfcost}(\fnistart(i)) = \bot \\
\end{cases}\) \\
Fixed Cost (activity granularity) & \dcost & \gact & \(a \in A\) & \(\indfc\utanyvariation(a)\) 
\newline \newline \(\utanyvariationnrm \in \{\utsinglenrm, \utsummationnrm\}\) & \(\sum _{i \in \fninst(a)} \indfc\utanyvariation(i)\)
\newline \newline The $\utanyvariationnrm$ indicates that the PPI can take multiple forms, in this case, $\utanyvariationnrm \in \{\utsinglenrm, \utsummationnrm\}$. If $\utanyvariationnrm = \utsinglenrm$, then the function considers single events of activity instances for cost calculations, i.e., it uses $\indfc\utsingle$; if $\utanyvariationnrm = \utsummationnrm$, then the function considers the sum of all events of activity instances for cost calculations, i.e., it uses $\indfc\utsummation$. \\
Fixed Cost (case granularity) & \dcost & \gcase & \(c \in C\) & \(\indfc\utanyvariation(c)\) 
\newline \newline \(\utanyvariationnrm \in \{\utsinglenrm, \utsummationnrm\}\) & \(\sum _{i \in \fninst(c)} \indfc\utanyvariation(i)\)
\newline \newline The $\utanyvariationnrm$ indicates that the PPI can take multiple forms, in this case, $\utanyvariationnrm \in \{\utsinglenrm, \utsummationnrm\}$. If $\utanyvariationnrm = \utsinglenrm$, then the function considers single events of activity instances for cost calculations, i.e., it uses $\indfc\utsingle$; if $\utanyvariationnrm = \utsummationnrm$, then the function considers the sum of all events of activity instances for cost calculations, i.e., it uses $\indfc\utsummation$. \\
Fixed Cost (group of cases granularity) & \dcost & \ggroup & \(g \subseteq C \colon\) \newline \(g \neq \emptyset\) & \(\indfc\utanyvariation(g)\) 
\newline \newline \(\utanyvariationnrm \in \{\utsinglenrm, \utsummationnrm\}\) & \(\sum _{c \in g} \indfc\utanyvariation(c)\)
\newline \newline The $\utanyvariationnrm$ indicates that the PPI can take multiple forms, in this case, $\utanyvariationnrm \in \{\utsinglenrm, \utsummationnrm\}$. If $\utanyvariationnrm = \utsinglenrm$, then the function considers single events of activity instances for cost calculations, i.e., it uses $\indfc\utsingle$; if $\utanyvariationnrm = \utsummationnrm$, then the function considers the sum of all events of activity instances for cost calculations, i.e., it uses $\indfc\utsummation$. \\
Expected Fixed Cost & \dcost & \ggroup & \(g \subseteq C \colon\) \newline \(g \neq \emptyset\) & \(\utexpected\indfc\utanyvariation(g)\) 
\newline \newline \(\utanyvariationnrm \in \{\utsinglenrm, \utsummationnrm\}\) & \(\frac{\indfc\utanyvariation(g)}{\indc(g)}\)
\newline \newline The $\utanyvariationnrm$ indicates that the PPI can take multiple forms, in this case, $\utanyvariationnrm \in \{\utsinglenrm, \utsummationnrm\}$. If $\utanyvariationnrm = \utsinglenrm$, then the function considers single events of activity instances for cost calculations, i.e., it uses $\indfc\utsingle$; if $\utanyvariationnrm = \utsummationnrm$, then the function considers the sum of all events of activity instances for cost calculations, i.e., it uses $\indfc\utsummation$. \\
Human Resource and Case count Ratio (group of cases granularity) & \dcost & \ggroup & \(g \subseteq C \colon\) \newline \(g \neq \emptyset\) & \(\indhrcr(g)\) & \(\frac{\indhr(g)}{\indc(g)}\) \\
Human Resource count (activity granularity) & \dcost & \gact & \(a \in A\) & \(\indhr(a)\) & \(|\fnhres(a)|\) \\
Human Resource count (case granularity) & \dcost & \gcase & \(c \in C\) & \(\indhr(c)\) & \(|\fnhres(c)|\) \\
Human Resource count (group of cases granularity) & \dcost & \ggroup & \(g \subseteq C \colon\) \newline \(g \neq \emptyset\) & \(\indhr(g)\) & \(|\cup_{c \in g} \fnhres(c)|\) \\
Expected Human Resource count & \dcost & \ggroup & \(g \subseteq C \colon\) \newline \(g \neq \emptyset\) & \(\utexpected\indhr(g)\) & \(\frac{\sum _{c \in g} \indhr(c)}{\indc(g)}\) \\
Inventory Cost considering single events of activity instances (activity instance granularity) & \dcost & \ginst & \(i \in I\) & \(\indic\utsingle(i)\) & \(\begin{cases}
\attribute{\aticost}(\fnicomplete(i)) & \utif \attribute{\aticost}(\fnicomplete(i)) \neq \bot \\
\attribute{\aticost}(\fnistart(i)) & \utif \attribute{\aticost}(\fnicomplete(i)) = \bot \land \attribute{\aticost}(\fnistart(i)) \neq \bot \\
\bot & \utif \attribute{\aticost}(\fnicomplete(i)) = \bot \land \attribute{\aticost}(\fnistart(i)) = \bot \\
\end{cases}\) \\
Inventory Cost considering the sum of all events of activity instances (activity instance granularity) & \dcost & \ginst & \(i \in I\) & \(\indic\utsummation(i)\) & \(\begin{cases}
\attribute{\aticost}(\fnistart(i)) + & \\ \attribute{\aticost}(\fnicomplete(i)) & \utif \attribute{\aticost}(\fnicomplete(i)) \neq \bot \land \attribute{\aticost}(\fnistart(i)) \neq \bot \\
\bot & \utif \attribute{\aticost}(\fnicomplete(i)) = \bot \lor \attribute{\aticost}(\fnistart(i)) = \bot \\
\end{cases}\) \\
Inventory Cost (activity granularity) & \dcost & \gact & \(a \in A\) & \(\indic\utanyvariation(a)\) 
\newline \newline \(\utanyvariationnrm \in \{\utsinglenrm, \utsummationnrm\}\) & \(\sum _{i \in \fninst(a)} \indic\utanyvariation(i)\)
\newline \newline The $\utanyvariationnrm$ indicates that the PPI can take multiple forms, in this case, $\utanyvariationnrm \in \{\utsinglenrm, \utsummationnrm\}$. If $\utanyvariationnrm = \utsinglenrm$, then the function considers single events of activity instances for cost calculations, i.e., it uses $\indic\utsingle$; if $\utanyvariationnrm = \utsummationnrm$, then the function considers the sum of all events of activity instances for cost calculations, i.e., it uses $\indic\utsummation$. \\
Inventory Cost (case granularity) & \dcost & \gcase & \(c \in C\) & \(\indic\utanyvariation(c)\) 
\newline \newline \(\utanyvariationnrm \in \{\utsinglenrm, \utsummationnrm\}\) & \(\sum _{i \in \fninst(c)} \indic\utanyvariation(i)\)
\newline \newline The $\utanyvariationnrm$ indicates that the PPI can take multiple forms, in this case, $\utanyvariationnrm \in \{\utsinglenrm, \utsummationnrm\}$. If $\utanyvariationnrm = \utsinglenrm$, then the function considers single events of activity instances for cost calculations, i.e., it uses $\indic\utsingle$; if $\utanyvariationnrm = \utsummationnrm$, then the function considers the sum of all events of activity instances for cost calculations, i.e., it uses $\indic\utsummation$. \\
Inventory Cost (group of cases granularity) & \dcost & \ggroup & \(g \subseteq C \colon\) \newline \(g \neq \emptyset\) & \(\indic\utanyvariation(g)\) 
\newline \newline \(\utanyvariationnrm \in \{\utsinglenrm, \utsummationnrm\}\) & \(\sum _{c \in g} \indic\utanyvariation(c)\)
\newline \newline The $\utanyvariationnrm$ indicates that the PPI can take multiple forms, in this case, $\utanyvariationnrm \in \{\utsinglenrm, \utsummationnrm\}$. If $\utanyvariationnrm = \utsinglenrm$, then the function considers single events of activity instances for cost calculations, i.e., it uses $\indic\utsingle$; if $\utanyvariationnrm = \utsummationnrm$, then the function considers the sum of all events of activity instances for cost calculations, i.e., it uses $\indic\utsummation$. \\
Expected Inventory Cost & \dcost & \ggroup & \(g \subseteq C \colon\) \newline \(g \neq \emptyset\) & \(\utexpected\indic\utanyvariation(g)\) 
\newline \newline \(\utanyvariationnrm \in \{\utsinglenrm, \utsummationnrm\}\) & \(\frac{\indic\utanyvariation(g)}{\indc(g)}\)
\newline \newline The $\utanyvariationnrm$ indicates that the PPI can take multiple forms, in this case, $\utanyvariationnrm \in \{\utsinglenrm, \utsummationnrm\}$. If $\utanyvariationnrm = \utsinglenrm$, then the function considers single events of activity instances for cost calculations, i.e., it uses $\indic\utsingle$; if $\utanyvariationnrm = \utsummationnrm$, then the function considers the sum of all events of activity instances for cost calculations, i.e., it uses $\indic\utsummation$. \\
Labor Cost and Total Cost Ratio (activity instance granularity) & \dcost & \ginst & \(i \in I\) & \(\indlctcr\utanyvariation(i)\) 
\newline \newline \(\utanyvariationnrm \in \{\utsinglenrm, \utsummationnrm\}\) & \(\frac{\indlc\utanyvariation(i)}{\indtc\utanyvariation(i)}\)
\newline \newline The $\utanyvariationnrm$ indicates that the PPI can take multiple forms, in this case, $\utanyvariationnrm \in \{\utsinglenrm, \utsummationnrm\}$. If $\utanyvariationnrm = \utsinglenrm$, then the function considers single events of activity instances for cost calculations, i.e., it uses $\indlc\utsingle$ and $\indtc\utsingle$; if $\utanyvariationnrm = \utsummationnrm$, then the function considers the sum of all events of activity instances for cost calculations, i.e., it uses $\indlc\utsummation$ and $\indtc\utsummation$. \\
Labor Cost and Total Cost Ratio (activity granularity) & \dcost & \gact & \(a \in A\) & \(\indlctcr\utanyvariation(a)\) 
\newline \newline \(\utanyvariationnrm \in \{\utsinglenrm, \utsummationnrm\}\) & \(\frac{\indlc\utanyvariation(a)}{\indtc\utanyvariation(a)}\)
\newline \newline The $\utanyvariationnrm$ indicates that the PPI can take multiple forms, in this case, $\utanyvariationnrm \in \{\utsinglenrm, \utsummationnrm\}$. If $\utanyvariationnrm = \utsinglenrm$, then the function considers single events of activity instances for cost calculations, i.e., it uses $\indlc\utsingle$ and $\indtc\utsingle$; if $\utanyvariationnrm = \utsummationnrm$, then the function considers the sum of all events of activity instances for cost calculations, i.e., it uses $\indlc\utsummation$ and $\indtc\utsummation$. \\
Labor Cost and Total Cost Ratio (case granularity) & \dcost & \gcase & \(c \in C\) & \(\indlctcr\utanyvariation(c)\) 
\newline \newline \(\utanyvariationnrm \in \{\utsinglenrm, \utsummationnrm\}\) & \(\frac{\indlc\utanyvariation(c)}{\indtc\utanyvariation(c)}\)
\newline \newline The $\utanyvariationnrm$ indicates that the PPI can take multiple forms, in this case, $\utanyvariationnrm \in \{\utsinglenrm, \utsummationnrm\}$. If $\utanyvariationnrm = \utsinglenrm$, then the function considers single events of activity instances for cost calculations, i.e., it uses $\indlc\utsingle$ and $\indtc\utsingle$; if $\utanyvariationnrm = \utsummationnrm$, then the function considers the sum of all events of activity instances for cost calculations, i.e., it uses $\indlc\utsummation$ and $\indtc\utsummation$. \\
Labor Cost and Total Cost Ratio (group of cases granularity) & \dcost & \ggroup & \(g \subseteq C \colon\) \newline \(g \neq \emptyset\) & \(\indlctcr\utanyvariation(g)\) 
\newline \newline \(\utanyvariationnrm \in \{\utsinglenrm, \utsummationnrm\}\) & \(\frac{\indlc\utanyvariation(g)}{\indtc\utanyvariation(g)}\)
\newline \newline The $\utanyvariationnrm$ indicates that the PPI can take multiple forms, in this case, $\utanyvariationnrm \in \{\utsinglenrm, \utsummationnrm\}$. If $\utanyvariationnrm = \utsinglenrm$, then the function considers single events of activity instances for cost calculations, i.e., it uses $\indlc\utsingle$ and $\indtc\utsingle$; if $\utanyvariationnrm = \utsummationnrm$, then the function considers the sum of all events of activity instances for cost calculations, i.e., it uses $\indlc\utsummation$ and $\indtc\utsummation$. \\
Expected Labor Cost and Total Cost Ratio & \dcost & \ggroup & \(g \subseteq C \colon\) \newline \(g \neq \emptyset\) & \(\utexpected\indlctcr\utanyvariation(g)\) 
\newline \newline \(\utanyvariationnrm \in \{\utsinglenrm, \utsummationnrm\}\) & \(\frac{\indlc\utanyvariation(g)}{\indtc\utanyvariation(g)}\)
\newline \newline The $\utanyvariationnrm$ indicates that the PPI can take multiple forms, in this case, $\utanyvariationnrm \in \{\utsinglenrm, \utsummationnrm\}$. If $\utanyvariationnrm = \utsinglenrm$, then the function considers single events of activity instances for cost calculations, i.e., it uses $\indlc\utsingle$ and $\indtc\utsingle$; if $\utanyvariationnrm = \utsummationnrm$, then the function considers the sum of all events of activity instances for cost calculations, i.e., it uses $\indlc\utsummation$ and $\indtc\utsummation$. \\
Labor Cost considering single events of activity instances (activity instance granularity) & \dcost & \ginst & \(i \in I\) & \(\indlc\utsingle(i)\) & \(\begin{cases}
\attribute{\atlcost}(\fnicomplete(i)) & \utif \attribute{\atlcost}(\fnicomplete(i)) \neq \bot \\
\attribute{\atlcost}(\fnistart(i)) & \utif \attribute{\atlcost}(\fnicomplete(i)) = \bot \land \attribute{\atlcost}(\fnistart(i)) \neq \bot \\
\bot & \utif \attribute{\atlcost}(\fnicomplete(i)) = \bot \land \attribute{\atlcost}(\fnistart(i)) = \bot \\
\end{cases}\) \\
Labor Cost considering the sum of all events of activity instances (activity instance granularity) & \dcost & \ginst & \(i \in I\) & \(\indlc\utsummation(i)\) & \(\begin{cases}
\attribute{\atlcost}(\fnistart(i)) + & \\ \attribute{\atlcost}(\fnicomplete(i)) & \utif \attribute{\atlcost}(\fnicomplete(i)) \neq \bot \land \attribute{\atlcost}(\fnistart(i)) \neq \bot \\
\bot & \utif \attribute{\atlcost}(\fnicomplete(i)) = \bot \lor \attribute{\atlcost}(\fnistart(i)) = \bot \\
\end{cases}\) \\
Labor Cost (activity granularity) & \dcost & \gact & \(a \in A\) & \(\indlc\utanyvariation(a)\) 
\newline \newline \(\utanyvariationnrm \in \{\utsinglenrm, \utsummationnrm\}\) & \(\sum _{i \in \fninst(a)} \indlc\utanyvariation(i)\)
\newline \newline The $\utanyvariationnrm$ indicates that the PPI can take multiple forms, in this case, $\utanyvariationnrm \in \{\utsinglenrm, \utsummationnrm\}$. If $\utanyvariationnrm = \utsinglenrm$, then the function considers single events of activity instances for cost calculations, i.e., it uses $\indlc\utsingle$; if $\utanyvariationnrm = \utsummationnrm$, then the function considers the sum of all events of activity instances for cost calculations, i.e., it uses $\indlc\utsummation$. \\
Labor Cost (case granularity) & \dcost & \gcase & \(c \in C\) & \(\indlc\utanyvariation(c)\) 
\newline \newline \(\utanyvariationnrm \in \{\utsinglenrm, \utsummationnrm\}\) & \(\sum _{i \in \fninst(c)} \indlc\utanyvariation(i)\)
\newline \newline The $\utanyvariationnrm$ indicates that the PPI can take multiple forms, in this case, $\utanyvariationnrm \in \{\utsinglenrm, \utsummationnrm\}$. If $\utanyvariationnrm = \utsinglenrm$, then the function considers single events of activity instances for cost calculations, i.e., it uses $\indlc\utsingle$; if $\utanyvariationnrm = \utsummationnrm$, then the function considers the sum of all events of activity instances for cost calculations, i.e., it uses $\indlc\utsummation$. \\
Labor Cost (group of cases granularity) & \dcost & \ggroup & \(g \subseteq C \colon\) \newline \(g \neq \emptyset\) & \(\indlc\utanyvariation(g)\) 
\newline \newline \(\utanyvariationnrm \in \{\utsinglenrm, \utsummationnrm\}\) & \(\sum _{c \in g} \indlc\utanyvariation(c)\)
\newline \newline The $\utanyvariationnrm$ indicates that the PPI can take multiple forms, in this case, $\utanyvariationnrm \in \{\utsinglenrm, \utsummationnrm\}$. If $\utanyvariationnrm = \utsinglenrm$, then the function considers single events of activity instances for cost calculations, i.e., it uses $\indlc\utsingle$; if $\utanyvariationnrm = \utsummationnrm$, then the function considers the sum of all events of activity instances for cost calculations, i.e., it uses $\indlc\utsummation$. \\
Expected Labor Cost & \dcost & \ggroup & \(g \subseteq C \colon\) \newline \(g \neq \emptyset\) & \(\utexpected\indlc\utanyvariation(g)\) 
\newline \newline \(\utanyvariationnrm \in \{\utsinglenrm, \utsummationnrm\}\) & \(\frac{\indlc\utanyvariation(g)}{\indc(g)}\)
\newline \newline The $\utanyvariationnrm$ indicates that the PPI can take multiple forms, in this case, $\utanyvariationnrm \in \{\utsinglenrm, \utsummationnrm\}$. If $\utanyvariationnrm = \utsinglenrm$, then the function considers single events of activity instances for cost calculations, i.e., it uses $\indlc\utsingle$; if $\utanyvariationnrm = \utsummationnrm$, then the function considers the sum of all events of activity instances for cost calculations, i.e., it uses $\indlc\utsummation$. \\
Maintenance Cost (case granularity) & \dcost & \gcase & \(c \in C\) & \(\indmainc(c)\) & \(\attribute{\atmaincost}(c)\) \\
Maintenance Cost (group of cases granularity) & \dcost & \ggroup & \(g \subseteq C \colon\) \newline \(g \neq \emptyset\) & \(\indmainc(g)\) & \(\sum _{c \in g} \indmainc(c)\) \\
Expected Maintenance Cost & \dcost & \ggroup & \(g \subseteq C \colon\) \newline \(g \neq \emptyset\) & \(\utexpected\indmainc(g)\) & \(\frac{\indmainc(g)}{\indc(g)}\) \\
Missed Deadline Cost (case granularity) & \dcost & \gcase & \(c \in C\) & \(\indmdc(c)\) & \(\attribute{\atmdcost}(c)\) \\
Missed Deadline Cost (group of cases granularity) & \dcost & \ggroup & \(g \subseteq C \colon\) \newline \(g \neq \emptyset\) & \(\indmdc(g)\) & \(\sum _{c \in g} \indmdc(c)\) \\
Expected Missed Deadline Cost & \dcost & \ggroup & \(g \subseteq C \colon\) \newline \(g \neq \emptyset\) & \(\utexpected\indmdc(g)\) & \(\frac{\indmdc(g)}{\indc(g)}\) \\
Overhead Cost (case granularity) & \dcost & \gcase & \(c \in C\) & \(\indoc\utanyvariation(c)\) 
\newline \newline \(\utanyvariationnrm \in \{\utsinglenrm, \utsummationnrm\}\) & \(\sum _{i \in \fninst(c), \fnact(i) \notin \lsdcl} \indtc\utanyvariation(i)\)
\newline \newline The $\utanyvariationnrm$ indicates that the PPI can take multiple forms, in this case, $\utanyvariationnrm \in \{\utsinglenrm, \utsummationnrm\}$. If $\utanyvariationnrm = \utsinglenrm$, then the function considers single events of activity instances for cost calculations, i.e., it uses $\indtc\utsingle$; if $\utanyvariationnrm = \utsummationnrm$, then the function considers the sum of all events of activity instances for cost calculations, i.e., it uses $\indtc\utsummation$. \\
Overhead Cost (group of cases granularity) & \dcost & \ggroup & \(g \subseteq C \colon\) \newline \(g \neq \emptyset\) & \(\indoc\utanyvariation(g)\) 
\newline \newline \(\utanyvariationnrm \in \{\utsinglenrm, \utsummationnrm\}\) & \(\sum _{c \in g} \indoc\utanyvariation(c)\)
\newline \newline The $\utanyvariationnrm$ indicates that the PPI can take multiple forms, in this case, $\utanyvariationnrm \in \{\utsinglenrm, \utsummationnrm\}$. If $\utanyvariationnrm = \utsinglenrm$, then the function considers single events of activity instances for cost calculations, i.e., it uses $\indoc\utsingle$; if $\utanyvariationnrm = \utsummationnrm$, then the function considers the sum of all events of activity instances for cost calculations, i.e., it uses $\indoc\utsummation$. \\
Expected Overhead Cost & \dcost & \ggroup & \(g \subseteq C \colon\) \newline \(g \neq \emptyset\) & \(\utexpected\indoc\utanyvariation(g)\) 
\newline \newline \(\utanyvariationnrm \in \{\utsinglenrm, \utsummationnrm\}\) & \(\frac{\indoc\utanyvariation(g)}{\indc(g)}\)
\newline \newline The $\utanyvariationnrm$ indicates that the PPI can take multiple forms, in this case, $\utanyvariationnrm \in \{\utsinglenrm, \utsummationnrm\}$. If $\utanyvariationnrm = \utsinglenrm$, then the function considers single events of activity instances for cost calculations, i.e., it uses $\indoc\utsingle$; if $\utanyvariationnrm = \utsummationnrm$, then the function considers the sum of all events of activity instances for cost calculations, i.e., it uses $\indoc\utsummation$. \\
Resource count (activity granularity) & \dcost & \gact & \(a \in A\) & \(\indr(a)\) & \(|\fnres(a)|\) \\
Resource count (case granularity) & \dcost & \gcase & \(c \in C\) & \(\indr(c)\) & \(|\fnres(c)|\) \\
Resource count (group of cases granularity) & \dcost & \ggroup & \(g \subseteq C \colon\) \newline \(g \neq \emptyset\) & \(\indr(g)\) & \(|\cup_{c \in g} \fnres(c)|\) \\
Expected Resource count & \dcost & \ggroup & \(g \subseteq C \colon\) \newline \(g \neq \emptyset\) & \(\utexpected\indr(g)\) & \(\frac{\sum _{c \in g} \indr(c)}{\indc(g)}\) \\
Rework Cost (activity granularity) & \dcost & \gact & \(a \in A\) & \(\indrc\utanyvariation(a)\)
\newline \newline \(\utanyvariationnrm \in \{\utsinglenrm, \utsummationnrm\}\) & \(\indtc(a) - \sum _{c \in C} \fnfirstinsttc\utanyvariation(c, a)\)
\newline \newline The $\utanyvariationnrm$ indicates that the PPI can take multiple forms, in this case, $\utanyvariationnrm \in \{\utsinglenrm, \utsummationnrm\}$. If $\utanyvariationnrm = \utsinglenrm$, then the function considers single events of activity instances for cost calculations, i.e., it uses $\fnfirstinsttc\utsingle$; if $\utanyvariationnrm = \utsummationnrm$, then the function considers the sum of all events of activity instances for cost calculations, i.e., it uses $\fnfirstinsttc\utsummation$. \\
Rework Cost (case granularity) & \dcost & \gcase & \(c \in C\) & \(\indrc\utanyvariation(c)\)
\newline \newline \(\utanyvariationnrm \in \{\utsinglenrm, \utsummationnrm\}\) & \(\indtc(c) - \sum _{a \in A} \fnfirstinsttc\utanyvariation(c, a)\)
\newline \newline The $\utanyvariationnrm$ indicates that the PPI can take multiple forms, in this case, $\utanyvariationnrm \in \{\utsinglenrm, \utsummationnrm\}$. If $\utanyvariationnrm = \utsinglenrm$, then the function considers single events of activity instances for cost calculations, i.e., it uses $\fnfirstinsttc\utsingle$; if $\utanyvariationnrm = \utsummationnrm$, then the function considers the sum of all events of activity instances for cost calculations, i.e., it uses $\fnfirstinsttc\utsummation$. \\
Rework Cost (group of cases granularity) & \dcost & \ggroup & \(g \subseteq C \colon\) \newline \(g \neq \emptyset\) & \(\indrc\utanyvariation(g)\)
\newline \newline \(\utanyvariationnrm \in \{\utsinglenrm, \utsummationnrm\}\) & \(\sum _{c \in g} \indrc\utanyvariation(c)\)
\newline \newline The $\utanyvariationnrm$ indicates that the PPI can take multiple forms, in this case, $\utanyvariationnrm \in \{\utsinglenrm, \utsummationnrm\}$. If $\utanyvariationnrm = \utsinglenrm$, then the function considers single events of activity instances for cost calculations, i.e., it uses $\indrc\utsingle$; if $\utanyvariationnrm = \utsummationnrm$, then the function considers the sum of all events of activity instances for cost calculations, i.e., it uses $\indrc\utsummation$. \\
Expected Rework Cost & \dcost & \ggroup & \(g \subseteq C \colon\) \newline \(g \neq \emptyset\) & \(\utexpected\indrc\utanyvariation(g)\)
\newline \newline \(\utanyvariationnrm \in \{\utsinglenrm, \utsummationnrm\}\) & \(\frac{\indrc\utanyvariation(g)}{\indc(g)}\)
\newline \newline The $\utanyvariationnrm$ indicates that the PPI can take multiple forms, in this case, $\utanyvariationnrm \in \{\utsinglenrm, \utsummationnrm\}$. If $\utanyvariationnrm = \utsinglenrm$, then the function considers single events of activity instances for cost calculations, i.e., it uses $\indrc\utsingle$; if $\utanyvariationnrm = \utsummationnrm$, then the function considers the sum of all events of activity instances for cost calculations, i.e., it uses $\indrc\utsummation$. \\
Rework Count (activity granularity) & \dcost & \gact & \(a \in A\) & \(\indrewc(a)\) & \(\sum _{c \in C} \fnmax(0, \fncount(c, a) - 1)\) \\
Rework Count (case granularity) & \dcost & \gcase & \(c \in C\) & \(\indrewc(c)\) & \(\sum _{a \in A} \fnmax(0, \fncount(c, a) - 1)\) \\
Rework Count (group of cases granularity) & \dcost & \ggroup & \(g \subseteq C \colon\) \newline \(g \neq \emptyset\) & \(\indrewc(g)\) & \(\sum _{c \in g} \indrewc(c)\) \\
Expected Rework Count & \dcost & \ggroup & \(g \subseteq C \colon\) \newline \(g \neq \emptyset\) & \(\utexpected\indrewc(g)\) & \(\frac{\indrewc(g)}{\indc(g)}\) \\
Rework Percentage (activity granularity) & \dcost & \gact & \(a \in A\) & \(\indrewp(a)\) & \(\frac{\indrewc(a)}{\indi(a)}\) \\
Rework Percentage (case granularity) & \dcost & \gcase & \(c \in C\) & \(\indrewp(c)\) & \(\frac{\indrewc(c)}{\indi(c)}\) \\
Rework Percentage (group of cases granularity) & \dcost & \ggroup & \(g \subseteq C \colon\) \newline \(g \neq \emptyset\) & \(\indrewp(g)\) & \(\frac{\indrewc(g)}{\indi(g)}\) \\
Expected Rework Percentage & \dcost & \ggroup & \(g \subseteq C \colon\) \newline \(g \neq \emptyset\) & \(\utexpected\indrewp(g)\) & \(\frac{\sum _{c \in g} \indrewp(c)}{\indc(g)}\) \\
Total Cost considering single events of activity instances (activity instance granularity) & \dcost & \ginst & \(i \in I\) & \(\indtc\utsingle(i)\) & \(\begin{cases}
\attribute{\attcost}(\fnicomplete(i)) & \utif \attribute{\attcost}(\fnicomplete(i)) \neq \bot \\
\attribute{\attcost}(\fnistart(i)) & \utif \attribute{\attcost}(\fnicomplete(i)) = \bot \land \attribute{\attcost}(\fnistart(i)) \neq \bot \\
\bot & \utif \attribute{\attcost}(\fnicomplete(i)) = \bot \land \attribute{\attcost}(\fnistart(i)) = \bot \\
\end{cases}\) \\
Total Cost considering the sum of all events of activity instances (activity instance granularity) & \dcost & \ginst & \(i \in I\) & \(\indtc\utsummation(i)\) & \(\begin{cases}
\attribute{\attcost}(\fnistart(i)) + & \\ \attribute{\attcost}(\fnicomplete(i)) & \utif \attribute{\attcost}(\fnicomplete(i)) \neq \bot \land \attribute{\attcost}(\fnistart(i)) \neq \bot \\
\bot & \utif \attribute{\attcost}(\fnicomplete(i)) = \bot \lor \attribute{\attcost}(\fnistart(i)) = \bot \\
\end{cases}\) \\
Total Cost (activity granularity) & \dcost & \gact & \(a \in A\) & \(\indtc\utanyvariation(a)\)
\newline \newline \(\utanyvariationnrm \in \{\utsinglenrm, \utsummationnrm\}\) & \(\sum _{i \in \fninst(a)} \indtc\utanyvariation(i)\)
\newline \newline The $\utanyvariationnrm$ indicates that the PPI can take multiple forms, in this case, $\utanyvariationnrm \in \{\utsinglenrm, \utsummationnrm\}$. If $\utanyvariationnrm = \utsinglenrm$, then the function considers single events of activity instances for cost calculations, i.e., it uses $\indtc\utsingle$; if $\utanyvariationnrm = \utsummationnrm$, then the function considers the sum of all events of activity instances for cost calculations, i.e., it uses $\indtc\utsummation$. \\
Total Cost (case granularity) & \dcost & \gcase & \(c \in C\) & \(\indtc\utanyvariation(c)\)
\newline \newline \(\utanyvariationnrm \in \{\utsinglenrm, \utsummationnrm\}\) & \(\sum _{i \in \fninst(c)} \indtc\utanyvariation(i)\)
\newline \newline The $\utanyvariationnrm$ indicates that the PPI can take multiple forms, in this case, $\utanyvariationnrm \in \{\utsinglenrm, \utsummationnrm\}$. If $\utanyvariationnrm = \utsinglenrm$, then the function considers single events of activity instances for cost calculations, i.e., it uses $\indtc\utsingle$; if $\utanyvariationnrm = \utsummationnrm$, then the function considers the sum of all events of activity instances for cost calculations, i.e., it uses $\indtc\utsummation$. \\
Total Cost (group of cases granularity) & \dcost & \ggroup & \(g \subseteq C \colon\) \newline \(g \neq \emptyset\) & \(\indtc\utanyvariation(g)\)
\newline \newline \(\utanyvariationnrm \in \{\utsinglenrm, \utsummationnrm\}\) & \(\sum _{c \in g} \indtc\utanyvariation(c)\)
\newline \newline The $\utanyvariationnrm$ indicates that the PPI can take multiple forms, in this case, $\utanyvariationnrm \in \{\utsinglenrm, \utsummationnrm\}$. If $\utanyvariationnrm = \utsinglenrm$, then the function considers single events of activity instances for cost calculations, i.e., it uses $\indtc\utsingle$; if $\utanyvariationnrm = \utsummationnrm$, then the function considers the sum of all events of activity instances for cost calculations, i.e., it uses $\indtc\utsummation$. \\
Expected Total Cost & \dcost & \ggroup & \(g \subseteq C \colon\) \newline \(g \neq \emptyset\) & \(\utexpected\indtc\utanyvariation(g)\)
\newline \newline \(\utanyvariationnrm \in \{\utsinglenrm, \utsummationnrm\}\) & \(\frac{\indtc\utanyvariation(g)}{\indc(g)}\)
\newline \newline The $\utanyvariationnrm$ indicates that the PPI can take multiple forms, in this case, $\utanyvariationnrm \in \{\utsinglenrm, \utsummationnrm\}$. If $\utanyvariationnrm = \utsinglenrm$, then the function considers single events of activity instances for cost calculations, i.e., it uses $\indtc\utsingle$; if $\utanyvariationnrm = \utsummationnrm$, then the function considers the sum of all events of activity instances for cost calculations, i.e., it uses $\indtc\utsummation$. \\
Total Cost and Lead Time Ratio (activity instance granularity) & \dcost & \ginst & \(i \in I\) & \(\indtcltr\utanyvariation(i)\)
\newline \newline \(\utanyvariationnrm \in \{\utsinglenrm, \utsummationnrm\}\) & \(\frac{\indtc\utanyvariation(i)}{\indlt(i)}\)
\newline \newline The $\utanyvariationnrm$ indicates that the PPI can take multiple forms, in this case, $\utanyvariationnrm \in \{\utsinglenrm, \utsummationnrm\}$. If $\utanyvariationnrm = \utsinglenrm$, then the function considers single events of activity instances for cost calculations, i.e., it uses $\indtc\utsingle$; if $\utanyvariationnrm = \utsummationnrm$, then the function considers the sum of all events of activity instances for cost calculations, i.e., it uses $\indtc\utsummation$. \\
Total Cost and Lead Time Ratio (activity granularity) & \dcost & \gact & \(a \in A\) & \(\indtcltr\utanyvariation(a)\)
\newline \newline \(\utanyvariationnrm \in \{\utsinglenrm, \utsummationnrm\}\) & \(\frac{\indtc\utanyvariation(a)}{\indlt(a)}\)
\newline \newline The $\utanyvariationnrm$ indicates that the PPI can take multiple forms, in this case, $\utanyvariationnrm \in \{\utsinglenrm, \utsummationnrm\}$. If $\utanyvariationnrm = \utsinglenrm$, then the function considers single events of activity instances for cost calculations, i.e., it uses $\indtc\utsingle$; if $\utanyvariationnrm = \utsummationnrm$, then the function considers the sum of all events of activity instances for cost calculations, i.e., it uses $\indtc\utsummation$. \\
Total Cost and Lead Time Ratio (case granularity) & \dcost & \gcase & \(c \in C\) & \(\indtcltr\utanyvariation(c)\)
\newline \newline \(\utanyvariationnrm \in \{\utsinglenrm, \utsummationnrm\}\) & \(\frac{\indtc\utanyvariation(c)}{\indlt(c)}\)
\newline \newline The $\utanyvariationnrm$ indicates that the PPI can take multiple forms, in this case, $\utanyvariationnrm \in \{\utsinglenrm, \utsummationnrm\}$. If $\utanyvariationnrm = \utsinglenrm$, then the function considers single events of activity instances for cost calculations, i.e., it uses $\indtc\utsingle$; if $\utanyvariationnrm = \utsummationnrm$, then the function considers the sum of all events of activity instances for cost calculations, i.e., it uses $\indtc\utsummation$. \\
Total Cost and Lead Time Ratio (group of cases granularity) & \dcost & \ggroup & \(g \subseteq C \colon\) \newline \(g \neq \emptyset\) & \(\indtcltr\utanyvariation(g)\)
\newline \newline \(\utanyvariationnrm \in \{\utsinglenrm, \utsummationnrm\}\) & \(\frac{\indtc\utanyvariation(g)}{\indlt(g)}\)
\newline \newline The $\utanyvariationnrm$ indicates that the PPI can take multiple forms, in this case, $\utanyvariationnrm \in \{\utsinglenrm, \utsummationnrm\}$. If $\utanyvariationnrm = \utsinglenrm$, then the function considers single events of activity instances for cost calculations, i.e., it uses $\indtc\utsingle$; if $\utanyvariationnrm = \utsummationnrm$, then the function considers the sum of all events of activity instances for cost calculations, i.e., it uses $\indtc\utsummation$. \\
Expected Total Cost and Lead Time Ratio & \dcost & \ggroup & \(g \subseteq C \colon\) \newline \(g \neq \emptyset\) & \(\utexpected\indtcltr\utanyvariation(g)\)
\newline \newline \(\utanyvariationnrm \in \{\utsinglenrm, \utsummationnrm\}\) & \(\frac{\indtc\utanyvariation(g)}{\sum _{c \in g} \indlt(c)}\)
\newline \newline The $\utanyvariationnrm$ indicates that the PPI can take multiple forms, in this case, $\utanyvariationnrm \in \{\utsinglenrm, \utsummationnrm\}$. If $\utanyvariationnrm = \utsinglenrm$, then the function considers single events of activity instances for cost calculations, i.e., it uses $\indtc\utsingle$; if $\utanyvariationnrm = \utsummationnrm$, then the function considers the sum of all events of activity instances for cost calculations, i.e., it uses $\indtc\utsummation$. \\
Total Cost and outcome Unit Ratio (activity instance granularity) & \dcost & \ginst & \(i \in I\) & \(\indtcur\utanyvariation(i)\)
\newline \newline \(\utanyvariationnrm \in \{\utsinglenrm, \utsummationnrm\}\) & \(\frac{\indtc\utanyvariation(i)}{\indu\utanyvariation(i)}\)
\newline \newline The $\utanyvariationnrm$ indicates that the PPI can take multiple forms, in this case, $\utanyvariationnrm \in \{\utsinglenrm, \utsummationnrm\}$. If $\utanyvariationnrm = \utsinglenrm$, then the function considers single events of activity instances for cost and outcome unit calculations, i.e., it uses $\indtc\utsingle$ and $\indu\utsingle$; if $\utanyvariationnrm = \utsummationnrm$, then the function considers the sum of all events of activity instances for cost and outcome unit calculations, i.e., it uses $\indtc\utsummation$ and $\indu\utsummation$. \\
Total Cost and outcome Unit Ratio (activity granularity) & \dcost & \gact & \(a \in A\) & \(\indtcur\utanyvariation(a)\)
\newline \newline \(\utanyvariationnrm \in \{\utsinglenrm, \utsummationnrm\}\) & \(\frac{\indtc\utanyvariation(a)}{\indu\utanyvariation(a)}\)
\newline \newline The $\utanyvariationnrm$ indicates that the PPI can take multiple forms, in this case, $\utanyvariationnrm \in \{\utsinglenrm, \utsummationnrm\}$. If $\utanyvariationnrm = \utsinglenrm$, then the function considers single events of activity instances for cost and outcome unit calculations, i.e., it uses $\indtc\utsingle$ and $\indu\utsingle$; if $\utanyvariationnrm = \utsummationnrm$, then the function considers the sum of all events of activity instances for cost and outcome unit calculations, i.e., it uses $\indtc\utsummation$ and $\indu\utsummation$. \\
Total Cost and outcome Unit Ratio (case granularity) & \dcost & \gcase & \(c \in C\) & \(\indtcur\utanyvariation(c)\)
\newline \newline \(\utanyvariationnrm \in \{\utsinglenrm, \utsummationnrm\}\) & \(\frac{\indtc\utanyvariation(c)}{\indu\utanyvariation(c)}\)
\newline \newline The $\utanyvariationnrm$ indicates that the PPI can take multiple forms, in this case, $\utanyvariationnrm \in \{\utsinglenrm, \utsummationnrm\}$. If $\utanyvariationnrm = \utsinglenrm$, then the function considers single events of activity instances for cost and outcome unit calculations, i.e., it uses $\indtc\utsingle$ and $\indu\utsingle$; if $\utanyvariationnrm = \utsummationnrm$, then the function considers the sum of all events of activity instances for cost and outcome unit calculations, i.e., it uses $\indtc\utsummation$ and $\indu\utsummation$. \\
Total Cost and outcome Unit Ratio (group of cases granularity) & \dcost & \ggroup & \(g \subseteq C \colon\) \newline \(g \neq \emptyset\) & \(\indtcur\utanyvariation(g)\)
\newline \newline \(\utanyvariationnrm \in \{\utsinglenrm, \utsummationnrm\}\) & \(\frac{\indtc\utanyvariation(g)}{\indu\utanyvariation(g)}\)
\newline \newline The $\utanyvariationnrm$ indicates that the PPI can take multiple forms, in this case, $\utanyvariationnrm \in \{\utsinglenrm, \utsummationnrm\}$. If $\utanyvariationnrm = \utsinglenrm$, then the function considers single events of activity instances for cost and outcome unit calculations, i.e., it uses $\indtc\utsingle$ and $\indu\utsingle$; if $\utanyvariationnrm = \utsummationnrm$, then the function considers the sum of all events of activity instances for cost and outcome unit calculations, i.e., it uses $\indtc\utsummation$ and $\indu\utsummation$. \\
Expected Total Cost and outcome Unit Ratio & \dcost & \ggroup & \(g \subseteq C \colon\) \newline \(g \neq \emptyset\) & \(\utexpected\indtcur\utanyvariation(g)\)
\newline \newline \(\utanyvariationnrm \in \{\utsinglenrm, \utsummationnrm\}\) & \(\frac{\indtc\utanyvariation(g)}{\indu\utanyvariation(g)}\)
\newline \newline The $\utanyvariationnrm$ indicates that the PPI can take multiple forms, in this case, $\utanyvariationnrm \in \{\utsinglenrm, \utsummationnrm\}$. If $\utanyvariationnrm = \utsinglenrm$, then the function considers single events of activity instances for cost and outcome unit calculations, i.e., it uses $\indtc\utsingle$ and $\indu\utsingle$; if $\utanyvariationnrm = \utsummationnrm$, then the function considers the sum of all events of activity instances for cost and outcome unit calculations, i.e., it uses $\indtc\utsummation$ and $\indu\utsummation$. \\
Total Cost and Service Time Ratio (activity instance granularity) & \dcost & \ginst & \(i \in I\) & \(\indtcstr\utanyvariation(i)\)
\newline \newline \(\utanyvariationnrm \in \{\utsinglenrm, \utsummationnrm\}\) & \(\frac{\indtc\utanyvariation(i)}{\indst(i)}\)
\newline \newline The $\utanyvariationnrm$ indicates that the PPI can take multiple forms, in this case, $\utanyvariationnrm \in \{\utsinglenrm, \utsummationnrm\}$. If $\utanyvariationnrm = \utsinglenrm$, then the function considers single events of activity instances for cost calculations, i.e., it uses $\indtc\utsingle$; if $\utanyvariationnrm = \utsummationnrm$, then the function considers the sum of all events of activity instances for cost calculations, i.e., it uses $\indtc\utsummation$. \\
Total Cost and Service Time Ratio (activity granularity) & \dcost & \gact & \(a \in A\) & \(\indtcstr\utanyvariation(a)\)
\newline \newline \(\utanyvariationnrm \in \{\utsinglenrm, \utsummationnrm\}\) & \(\frac{\indtc\utanyvariation(a)}{\indst(a)}\)
\newline \newline The $\utanyvariationnrm$ indicates that the PPI can take multiple forms, in this case, $\utanyvariationnrm \in \{\utsinglenrm, \utsummationnrm\}$. If $\utanyvariationnrm = \utsinglenrm$, then the function considers single events of activity instances for cost calculations, i.e., it uses $\indtc\utsingle$; if $\utanyvariationnrm = \utsummationnrm$, then the function considers the sum of all events of activity instances for cost calculations, i.e., it uses $\indtc\utsummation$. \\
Total Cost and Service Time Ratio (case granularity) & \dcost & \gcase & \(c \in C\) & \(\indtcstr\utanyvariation(c)\)
\newline \newline \(\utanyvariationnrm \in \{\utsinglenrm, \utsummationnrm\}\) & \(\frac{\indtc\utanyvariation(c)}{\indst(c)}\)
\newline \newline The $\utanyvariationnrm$ indicates that the PPI can take multiple forms, in this case, $\utanyvariationnrm \in \{\utsinglenrm, \utsummationnrm\}$. If $\utanyvariationnrm = \utsinglenrm$, then the function considers single events of activity instances for cost calculations, i.e., it uses $\indtc\utsingle$; if $\utanyvariationnrm = \utsummationnrm$, then the function considers the sum of all events of activity instances for cost calculations, i.e., it uses $\indtc\utsummation$. \\
Total Cost and Service Time Ratio (group of cases granularity) & \dcost & \ggroup & \(g \subseteq C \colon\) \newline \(g \neq \emptyset\) & \(\indtcstr\utanyvariation(g)\)
\newline \newline \(\utanyvariationnrm \in \{\utsinglenrm, \utsummationnrm\}\) & \(\frac{\indtc\utanyvariation(g)}{\indst(g)}\)
\newline \newline The $\utanyvariationnrm$ indicates that the PPI can take multiple forms, in this case, $\utanyvariationnrm \in \{\utsinglenrm, \utsummationnrm\}$. If $\utanyvariationnrm = \utsinglenrm$, then the function considers single events of activity instances for cost calculations, i.e., it uses $\indtc\utsingle$; if $\utanyvariationnrm = \utsummationnrm$, then the function considers the sum of all events of activity instances for cost calculations, i.e., it uses $\indtc\utsummation$. \\
Expected Total Cost and Service Time Ratio & \dcost & \ggroup & \(g \subseteq C \colon\) \newline \(g \neq \emptyset\) & \(\utexpected\indtcstr\utanyvariation(g)\)
\newline \newline \(\utanyvariationnrm \in \{\utsinglenrm, \utsummationnrm\}\) & \(\frac{\indtc\utanyvariation(g)}{\indst(g)}\)
\newline \newline The $\utanyvariationnrm$ indicates that the PPI can take multiple forms, in this case, $\utanyvariationnrm \in \{\utsinglenrm, \utsummationnrm\}$. If $\utanyvariationnrm = \utsinglenrm$, then the function considers single events of activity instances for cost calculations, i.e., it uses $\indtc\utsingle$; if $\utanyvariationnrm = \utsummationnrm$, then the function considers the sum of all events of activity instances for cost calculations, i.e., it uses $\indtc\utsummation$. \\
Transportation Cost (case granularity) & \dcost & \gcase & \(c \in C\) & \(\indtransc(c)\) & \(\attribute{\attranscost}(c)\) \\
Transportation Cost (group of cases granularity) & \dcost & \ggroup & \(g \subseteq C \colon\) \newline \(g \neq \emptyset\) & \(\indtransc(g)\) & \(\sum _{c \in g} \indtransc(c)\) \\
Expected Transportation Cost & \dcost & \ggroup & \(g \subseteq C \colon\) \newline \(g \neq \emptyset\) & \(\utexpected\indtransc(g)\) & \(\frac{\indtransc(g)}{\indc(g)}\) \\
Variable Cost considering single events of activity instances (activity instance granularity) & \dcost & \ginst & \(i \in I\) & \(\indvc\utsingle(i)\) & \(\begin{cases}
\attribute{\atvcost}(\fnicomplete(i)) & \utif \attribute{\atvcost}(\fnicomplete(i)) \neq \bot \\
\attribute{\atvcost}(\fnistart(i)) & \utif \attribute{\atvcost}(\fnicomplete(i)) = \bot \land \attribute{\atvcost}(\fnistart(i)) \neq \bot \\
\bot & \utif \attribute{\atvcost}(\fnicomplete(i)) = \bot \land \attribute{\atvcost}(\fnistart(i)) = \bot \\
\end{cases}\) \\
Variable Cost considering the sum of all events of activity instances (activity instance granularity) & \dcost & \ginst & \(i \in I\) & \(\indvc\utsummation(i)\) & \(\begin{cases}
\attribute{\atvcost}(\fnistart(i)) + & \\ \attribute{\atvcost}(\fnicomplete(i)) & \utif \attribute{\atvcost}(\fnicomplete(i)) \neq \bot \land \attribute{\atvcost}(\fnistart(i)) \neq \bot \\
\bot & \utif \attribute{\atvcost}(\fnicomplete(i)) = \bot \lor \attribute{\atvcost}(\fnistart(i)) = \bot \\
\end{cases}\) \\
Variable Cost (activity granularity) & \dcost & \gact & \(a \in A\) & \(\indvc\utanyvariation(a)\)
\newline \newline \(\utanyvariationnrm \in \{\utsinglenrm, \utsummationnrm\}\) & \(\sum _{i \in \fninst(a)} \indvc\utanyvariation(i)\)
\newline \newline The $\utanyvariationnrm$ indicates that the PPI can take multiple forms, in this case, $\utanyvariationnrm \in \{\utsinglenrm, \utsummationnrm\}$. If $\utanyvariationnrm = \utsinglenrm$, then the function considers single events of activity instances for cost calculations, i.e., it uses $\indvc\utsingle$; if $\utanyvariationnrm = \utsummationnrm$, then the function considers the sum of all events of activity instances for cost calculations, i.e., it uses $\indvc\utsummation$. \\
Variable Cost (case granularity) & \dcost & \gcase & \(c \in C\) & \(\indvc\utanyvariation(c)\)
\newline \newline \(\utanyvariationnrm \in \{\utsinglenrm, \utsummationnrm\}\) & \(\sum _{i \in \fninst(c)} \indvc\utanyvariation(i)\)
\newline \newline The $\utanyvariationnrm$ indicates that the PPI can take multiple forms, in this case, $\utanyvariationnrm \in \{\utsinglenrm, \utsummationnrm\}$. If $\utanyvariationnrm = \utsinglenrm$, then the function considers single events of activity instances for cost calculations, i.e., it uses $\indvc\utsingle$; if $\utanyvariationnrm = \utsummationnrm$, then the function considers the sum of all events of activity instances for cost calculations, i.e., it uses $\indvc\utsummation$. \\
Variable Cost (group of cases granularity) & \dcost & \ggroup & \(g \subseteq C \colon\) \newline \(g \neq \emptyset\) & \(\indvc\utanyvariation(g)\)
\newline \newline \(\utanyvariationnrm \in \{\utsinglenrm, \utsummationnrm\}\) & \(\sum _{c \in g} \indvc\utanyvariation(c)\)
\newline \newline The $\utanyvariationnrm$ indicates that the PPI can take multiple forms, in this case, $\utanyvariationnrm \in \{\utsinglenrm, \utsummationnrm\}$. If $\utanyvariationnrm = \utsinglenrm$, then the function considers single events of activity instances for cost calculations, i.e., it uses $\indvc\utsingle$; if $\utanyvariationnrm = \utsummationnrm$, then the function considers the sum of all events of activity instances for cost calculations, i.e., it uses $\indvc\utsummation$. \\
Expected Variable Cost & \dcost & \ggroup & \(g \subseteq C \colon\) \newline \(g \neq \emptyset\) & \(\utexpected\indvc\utanyvariation(g)\)
\newline \newline \(\utanyvariationnrm \in \{\utsinglenrm, \utsummationnrm\}\) & \(\frac{\indvc\utanyvariation(g)}{\indc(g)}\)
\newline \newline The $\utanyvariationnrm$ indicates that the PPI can take multiple forms, in this case, $\utanyvariationnrm \in \{\utsinglenrm, \utsummationnrm\}$. If $\utanyvariationnrm = \utsinglenrm$, then the function considers single events of activity instances for cost calculations, i.e., it uses $\indvc\utsingle$; if $\utanyvariationnrm = \utsummationnrm$, then the function considers the sum of all events of activity instances for cost calculations, i.e., it uses $\indvc\utsummation$. \\
Warehousing Cost (case granularity) & \dcost & \gcase & \(c \in C\) & \(\indwarec(c)\) & \(\attribute{\atwarecost}(c)\) \\
Warehousing Cost (group of cases granularity) & \dcost & \ggroup & \(g \subseteq C \colon\) \newline \(g \neq \emptyset\) & \(\indwarec(g)\) & \(\sum _{c \in g} \indwarec(c)\) \\
Expected Warehousing Cost & \dcost & \ggroup & \(g \subseteq C \colon\) \newline \(g \neq \emptyset\) & \(\utexpected\indwarec(g)\) & \(\frac{\indwarec(g)}{\indc(g)}\) \\
\bottomrule
\label{PPICostDefinitions}
\end{longtable}


\begin{longtable}{>{\hspace{0pt}}m{0.14\linewidth}>{\hspace{0pt}}c>{\hspace{0pt}}c>{\hspace{0pt}}m{0.1\linewidth}>{\hspace{0pt}}m{0.19\linewidth}>{\hspace{0pt}}m{0.473\linewidth}}
\toprule
\textbf{Indicator} & \textbf{D.} & \textbf{G.} & \textbf{F. Par.} & \textbf{Function Name} & \textbf{Formal Definition}\\
\midrule \endfirsthead
\toprule
\textbf{Indicator} & \textbf{D.} & \textbf{G.} & \textbf{F. Par.} & \textbf{Function Name} & \textbf{Formal Definition}\\
\midrule \endhead
\bottomrule \endfoot
\captionsetup{font=small}
\caption{Quality dimension PPI definitions} \endlastfoot
activity Instance count by Human Resource (activity granularity) & \dquality & \gact & \(a \in A\) \newline \(\uthres \in \universe{\athres}\) & \(\indibyhr(a, \uthres)\) & \(|\{i \in \fninst(a) \utst \fnhres(i) = \uthres\}|\) \\
activity Instance count by Human Resource (case granularity) & \dquality & \gcase & \(c \in C\) \newline \(\uthres \in \universe{\athres}\) & \(\indibyhr(c, \uthres)\) & \(|\{i \in \fninst(c) \utst \fnhres(i) = \uthres\}|\) \\
activity Instance count by Human Resource (group of cases granularity) & \dquality & \ggroup & \(g \subseteq C \colon\) \newline \(g \neq \emptyset\) \newline \(\uthres \in \universe{\athres}\) & \(\indibyhr(g, \uthres)\) & \(\sum _{c \in g} \indibyhr(c, \uthres)\) \\
Expected activity Instance count by Human Resource & \dquality & \ggroup & \(g \subseteq C \colon\) \newline \(g \neq \emptyset\) \newline \(\uthres \in \universe{\athres}\) & \(\utexpected\indibyhr(g, \uthres)\) & \(\frac{\indibyhr(g, \uthres)}{\indc(g)}\) \\
activity Instance count by Role (case granularity) & \dquality & \gcase & \(c \in C\) \newline \(\utrole \in \universe{\atrole}\) & \(\indibyrole(c, \utrole)\) & \(|\{i \in \fninst(c) \utst \fnrole(i) = \utrole\}|\) \\
activity Instance count by Role (group of cases granularity) & \dquality & \ggroup & \(g \subseteq C \colon\) \newline \(g \neq \emptyset\) \newline \(\utrole \in \universe{\atrole}\) & \(\indibyrole(g, \utrole)\) & \(\sum _{c \in g} \indibyrole(c, \utrole)\) \\
Expected activity Instance count by Role & \dquality & \ggroup & \(g \subseteq C \colon\) \newline \(g \neq \emptyset\) \newline \(\utrole \in \universe{\atrole}\) & \(\utexpected\indibyrole(g, \utrole)\) & \(\frac{\indibyrole(g, \utrole)}{\indc(g)}\) \\
Automated Activity count (case granularity) & \dquality & \gcase & \(c \in C\) & \(\indauta(c)\) & \(|\lsautl \cap \fnact(c)|\) \\
Automated Activity count (group of cases granularity) & \dquality & \ggroup & \(g \subseteq C \colon\) \newline \(g \neq \emptyset\) & \(\indauta(g)\) & \(|\lsautl \cap (\cup_{c \in g} \fnact(c))|\) \\
Expected Automated Activity count & \dquality & \ggroup & \(g \subseteq C \colon\) \newline \(g \neq \emptyset\) & \(\utexpected\indauta(g)\) & \(\frac{\sum _{c \in g} \indauta(c)}{\indc(g)}\) \\
Automated activity Instance count (case granularity) & \dquality & \gcase & \(c \in C\) & \(\indauti(c)\) & \(|\{i \in \fninst(c) \utst \fnact(i) \in \lsautl\}|\) \\
Automated activity Instance count (group of cases granularity) & \dquality & \ggroup & \(g \subseteq C \colon\) \newline \(g \neq \emptyset\) & \(\indauti(g)\) & \(\sum _{c \in g} \indauti(c)\) \\
Expected Automated activity Instance count & \dquality & \ggroup & \(g \subseteq C \colon\) \newline \(g \neq \emptyset\) & \(\utexpected\indauti(g)\) & \(\frac{\indauti(g)}{\indc(g)}\) \\
Case and Client count Ratio (group of cases granularity) & \dquality & \ggroup & \(g \subseteq C \colon\) \newline \(g \neq \emptyset\) & \(\indcclir(g)\) & \(\frac{\indc(g)}{\indcli(g)}\) \\
Case Count where Activity after Time Frame (group of cases granularity) & \dquality & \ggroup & \(g \subseteq C \colon\) \newline \(g \neq \emptyset\) \newline \(\utacta \in A\) \newline \(\utetime \in \universe{\attime}\) & \(\indcaaftfc(g, \utacta, \utetime)\) & \(|\{c \in g \utst \exists i \in \fninst(c, \utacta)[\fnistarttime(i) \geq \utetime]\}|\) \\
Case Count where Activity before Time Frame (group of cases granularity) & \dquality & \ggroup & \(g \subseteq C \colon\) \newline \(g \neq \emptyset\) \newline \(\utacta \in A\) \newline \(\utstime \in \universe{\attime}\) & \(\indcabetfc(g, \utacta, \utstime)\) & \(|\{c \in g \utst \exists i \in \fninst(c, \utacta)[\fnistarttime(i) \leq \utstime]\}|\) \\
Case Count where Activity during Time Frame (group of cases granularity) & \dquality & \ggroup & \(g \subseteq C \colon\) \newline \(g \neq \emptyset\) \newline \(\utacta \in A\) \newline \(\utstime, \utetime \in \universe{\attime} \colon\) \newline \(\utstime < \utetime\) & \(\indcaintfc(g, \utacta, \utstime, \utetime)\) & \(|\{c \in g \utst \exists i \in \fninst(c, \utacta)[\utstime \leq \fnistarttime(i) \land \fnistarttime(i) \leq \utetime]\}|\) \\
Case Count where end activity is A (group of cases granularity) & \dquality & \ggroup & \(g \subseteq C \colon\) \newline \(g \neq \emptyset\) \newline \(\utacta \in A\) & \(\indcendac(g, \utacta)\) & \(|\{c \in g \utst \exists i \in \fnendinstance(c)[\fnact(i) = \utacta]\}|\) \\
Case Count where start activity is A (group of cases granularity) & \dquality & \ggroup & \(g \subseteq C \colon\) \newline \(g \neq \emptyset\) \newline \(\utacta \in A\) & \(\indcstartac(g, \utacta)\) & \(|\{c \in g \utst \exists i \in \fnstartinstance(c)[\fnact(i) = \utacta]\}|\) \\
Case Count with Rework (group of cases granularity) & \dquality & \ggroup & \(g \subseteq C \colon\) \newline \(g \neq \emptyset\) & \(\indcrewc(g)\) & \(|\{c \in g \utst \indrewc(c) > 0\}|\) \\
Case Percentage where Activity after Time Frame (group of cases granularity) & \dquality & \ggroup & \(g \subseteq C \colon\) \newline \(g \neq \emptyset\) \newline \(\utacta \in A\) \newline \(\utetime \in \universe{\attime}\) & \(\indcaaftfp(g, \utacta, \utetime)\) & \(\frac{\indcaaftfc(g, \utacta, \utetime)}{\indc(g)}\) \\
Case Percentage where Activity before Time Frame (group of cases granularity) & \dquality & \ggroup & \(g \subseteq C \colon\) \newline \(g \neq \emptyset\) \newline \(\utacta \in A\) \newline \(\utstime \in \universe{\attime}\) & \(\indcabetfp(g, \utacta, \utstime)\) & \(\frac{\indcabetfc(g, \utacta, \utstime)}{\indc(g)}\) \\
Case Percentage where Activity during Time Frame (group of cases granularity) & \dquality & \ggroup & \(g \subseteq C \colon\) \newline \(g \neq \emptyset\) \newline \(\utacta \in A\) \newline \(\utstime, \utetime \in \universe{\attime} \colon\) \newline \(\utstime < \utetime\) & \(\indcaintfp(g, \utacta, \utstime, \utetime)\) & \(\frac{\indcaintfc(g, \utacta, \utstime, \utetime)}{\indc(g)}\) \\
Case Percentage where End Activity is A (group of cases granularity) & \dquality & \ggroup & \(g \subseteq C \colon\) \newline \(g \neq \emptyset\) \newline \(\utacta \in A\) & \(\indcendap(g, \utacta)\) & \(\frac{\indcendac(g, \utacta)}{\indc(g)}\) \\
Case Percentage where Start Activity is A (group of cases granularity) & \dquality & \ggroup & \(g \subseteq C \colon\) \newline \(g \neq \emptyset\) \newline \(\utacta \in A\) & \(\indcstartap(g, \utacta)\) & \(\frac{\indcstartac(g, \utacta)}{\indc(g)}\) \\
Case Percentage with Missed Deadline (group of cases granularity) & \dquality & \ggroup & \(g \subseteq C \colon\) \newline \(g \neq \emptyset\) \newline \(\utval \in \universe{\attime}\) & \(\indcmdp(g, \utval)\) & \(\frac{|\{c \in g \utst \fnendtime(c) > \utval\}|}{\indc(g)}\) \\
Case Percentage with Rework (group of cases granularity) & \dquality & \ggroup & \(g \subseteq C \colon\) \newline \(g \neq \emptyset\) & \(\indcrewp(g)\) & \(\frac{\indcrewc(g)}{\indc(g)}\) \\
Client count and Total Cost Ratio (activity granularity) & \dquality & \gact & \(a \in A\) & \(\indclitcr\utanyvariation(a)\)
\newline \newline \(\utanyvariationnrm \in \{\utsinglenrm, \utsummationnrm\}\) & \(\frac{\indcli(a)}{\indtc\utanyvariation(a)}\)
\newline \newline The $\utanyvariationnrm$ indicates that the PPI can take multiple forms, in this case, $\utanyvariationnrm \in \{\utsinglenrm, \utsummationnrm\}$. If $\utanyvariationnrm = \utsinglenrm$, then the function considers single events of activity instances for cost calculations, i.e., it uses $\indtc\utsingle$; if $\utanyvariationnrm = \utsummationnrm$, then the function considers the sum of all events of activity instances for cost calculations, i.e., it uses $\indtc\utsummation$. \\
Client count and Total Cost Ratio (group of cases granularity) & \dquality & \ggroup & \(g \subseteq C \colon\) \newline \(g \neq \emptyset\) & \(\indclitcr\utanyvariation(g)\)
\newline \newline \(\utanyvariationnrm \in \{\utsinglenrm, \utsummationnrm\}\) & \(\frac{\indcli(g)}{\indtc\utanyvariation(g)}\)
\newline \newline The $\utanyvariationnrm$ indicates that the PPI can take multiple forms, in this case, $\utanyvariationnrm \in \{\utsinglenrm, \utsummationnrm\}$. If $\utanyvariationnrm = \utsinglenrm$, then the function considers single events of activity instances for cost calculations, i.e., it uses $\indtc\utsingle$; if $\utanyvariationnrm = \utsummationnrm$, then the function considers the sum of all events of activity instances for cost calculations, i.e., it uses $\indtc\utsummation$. \\
Expected Client count and Total Cost Ratio & \dquality & \ggroup & \(g \subseteq C \colon\) \newline \(g \neq \emptyset\) & \(\utexpected\indclitcr\utanyvariation(g)\)
\newline \newline \(\utanyvariationnrm \in \{\utsinglenrm, \utsummationnrm\}\) & \(\frac{\indcli(g)}{\indtc\utanyvariation(g)}\)
\newline \newline The $\utanyvariationnrm$ indicates that the PPI can take multiple forms, in this case, $\utanyvariationnrm \in \{\utsinglenrm, \utsummationnrm\}$. If $\utanyvariationnrm = \utsinglenrm$, then the function considers single events of activity instances for cost calculations, i.e., it uses $\indtc\utsingle$; if $\utanyvariationnrm = \utsummationnrm$, then the function considers the sum of all events of activity instances for cost calculations, i.e., it uses $\indtc\utsummation$. \\
Desired Activity count (case granularity) & \dquality & \gcase & \(c \in C\) & \(\inddac(c)\) & \(| \lsdesl \cap \fnact(c)|\) \\
Desired Activity count (group of cases granularity) & \dquality & \ggroup & \(g \subseteq C \colon\) \newline \(g \neq \emptyset\) & \(\inddac(g)\) & \(|\lsdesl \cap (\cup_{c \in g} \fnact(c))|\) \\
Expected Desired Activity count & \dquality & \ggroup & \(g \subseteq C \colon\) \newline \(g \neq \emptyset\) & \(\utexpected\inddac(g)\) & \(\frac{\sum _{c \in g} \inddac(c)}{\indc(g)}\) \\
Human Resource count (activity granularity) & \dquality & \gact & \(a \in A\) & \(\indhr(a)\) & \(|\fnhres(a)|\) \\
Human Resource count (case granularity) & \dquality & \gcase & \(c \in C\) & \(\indhr(c)\) & \(|\fnhres(c)|\) \\
Human Resource count (group of cases granularity) & \dquality & \ggroup & \(g \subseteq C \colon\) \newline \(g \neq \emptyset\) & \(\indhr(g)\) & \(|\cup_{c \in g} \fnhres(c)|\) \\
Expected Human Resource count & \dquality & \ggroup & \(g \subseteq C \colon\) \newline \(g \neq \emptyset\) & \(\utexpected\indhr(g)\) & \(\frac{\sum _{c \in g} \indhr(c)}{\indc(g)}\) \\
Non-Automated Activity count (case granularity) & \dquality & \gcase & \(c \in C\) & \(\indnauta(c)\) & \(|\fnact(c) \backslash \lsautl|\) \\
Non-Automated Activity count (group of cases granularity) & \dquality & \ggroup & \(g \subseteq C \colon\) \newline \(g \neq \emptyset\) & \(\indnauta(g)\) & \(|(\cup _{c \in g} \fnact(c)) \backslash \lsautl|\) \\
Expected Non-Automated Activity count (group of cases granularity) & \dquality & \ggroup & \(g \subseteq C \colon\) \newline \(g \neq \emptyset\) & \(\indnauta(g)\) & \(\frac{\sum _{c \in g} \indnauta(c)}{\indc(g)}\) \\
Non-Automated activity Instance count (case granularity) & \dquality & \gcase & \(c \in C\) & \(\indnauti(c)\) & \(|\{i \in \fninst(c) \utst \fnact(i) \notin \lsautl\}|\) \\
Non-Automated activity Instance count (group of cases granularity) & \dquality & \ggroup & \(g \subseteq C \colon\) \newline \(g \neq \emptyset\) & \(\indnauti(g)\) & \(\sum _{c \in g} \indnauti(c)\) \\
Expected Non-Automated activity Instance count & \dquality & \ggroup & \(g \subseteq C \colon\) \newline \(g \neq \emptyset\) & \(\utexpected\indnauti(g)\) & \(\frac{\indnauti(g)}{\indc(g)}\) \\
outcome Unit count considering single events of activity instances (activity instance granularity) & \dquality & \ginst & \(i \in I\) & \(\indu\utsingle(i)\) & \(\begin{cases}
\attribute{\atunit}(\fnicomplete(i)) & \utif \attribute{\atunit}(\fnicomplete(i)) \neq \bot \\
\attribute{\atunit}(\fnistart(i)) & \utif \attribute{\atunit}(\fnicomplete(i)) = \bot \land \\ & \attribute{\atunit}(\fnistart(i)) \neq \bot \\
\bot & \utif \attribute{\atunit}(\fnicomplete(i)) = \bot \land \\ & \attribute{\atunit}(\fnistart(i)) = \bot \\
\end{cases}\) \\
outcome Unit count considering the sum of all events of activity instances (activity instance granularity) & \dquality & \ginst & \(i \in I\) & \(\indu\utsummation(i)\) & \(\begin{cases}
\attribute{\atunit}(\fnistart(i)) + & \utif \attribute{\atunit}(\fnicomplete(i)) \neq \bot \land \\ \attribute{\atunit}(\fnicomplete(i)) & \attribute{\atunit}(\fnistart(i)) \neq \bot \\
\bot & \utif \attribute{\atunit}(\fnicomplete(i)) = \bot \lor \\ & \attribute{\atunit}(\fnistart(i)) = \bot \\
\end{cases}\) \\
outcome Unit count (activity granularity) & \dquality & \gact & \(a \in A\) & \(\indu\utanyvariation(a)\)
\newline \newline \(\utanyvariationnrm \in \{\utsinglenrm, \utsummationnrm\}\) & \(\sum _{i \in \fninst(a)} \indu\utanyvariation(i)\)
\newline \newline The $\utanyvariationnrm$ indicates that the PPI can take multiple forms, in this case, $\utanyvariationnrm \in \{\utsinglenrm, \utsummationnrm\}$. If $\utanyvariationnrm = \utsinglenrm$, then the function considers single events of activity instances for outcome unit calculations, i.e., it uses $\indu\utsingle$; if $\utanyvariationnrm = \utsummationnrm$, then the function considers the sum of all events of activity instances for outcome unit calculations, i.e., it uses $\indu\utsummation$. \\
outcome Unit count (case granularity) & \dquality & \gcase & \(c \in C\) & \(\indu\utanyvariation(c)\)
\newline \newline \(\utanyvariationnrm \in \{\utsinglenrm, \utsummationnrm\}\) & \(\sum _{e \in \fninst(c)} \indu\utanyvariation(i)\)
\newline \newline The $\utanyvariationnrm$ indicates that the PPI can take multiple forms, in this case, $\utanyvariationnrm \in \{\utsinglenrm, \utsummationnrm\}$. If $\utanyvariationnrm = \utsinglenrm$, then the function considers single events of activity instances for outcome unit calculations, i.e., it uses $\indu\utsingle$; if $\utanyvariationnrm = \utsummationnrm$, then the function considers the sum of all events of activity instances for outcome unit calculations, i.e., it uses $\indu\utsummation$. \\
outcome Unit count (group of cases granularity) & \dquality & \ggroup & \(g \subseteq C \colon\) \newline \(g \neq \emptyset\) & \(\indu\utanyvariation(g)\)
\newline \newline \(\utanyvariationnrm \in \{\utsinglenrm, \utsummationnrm\}\) & \(\sum _{c \in g} \indu\utanyvariation(c)\)
\newline \newline The $\utanyvariationnrm$ indicates that the PPI can take multiple forms, in this case, $\utanyvariationnrm \in \{\utsinglenrm, \utsummationnrm\}$. If $\utanyvariationnrm = \utsinglenrm$, then the function considers single events of activity instances for outcome unit calculations, i.e., it uses $\indu\utsingle$; if $\utanyvariationnrm = \utsummationnrm$, then the function considers the sum of all events of activity instances for outcome unit calculations, i.e., it uses $\indu\utsummation$. \\
Expected outcome Unit count & \dquality & \ggroup & \(g \subseteq C \colon\) \newline \(g \neq \emptyset\) & \(\utexpected\indu\utanyvariation(g)\)
\newline \newline \(\utanyvariationnrm \in \{\utsinglenrm, \utsummationnrm\}\) & \(\frac{\indu\utanyvariation(g)}{\indc(g)}\)
\newline \newline The $\utanyvariationnrm$ indicates that the PPI can take multiple forms, in this case, $\utanyvariationnrm \in \{\utsinglenrm, \utsummationnrm\}$. If $\utanyvariationnrm = \utsinglenrm$, then the function considers single events of activity instances for outcome unit calculations, i.e., it uses $\indu\utsingle$; if $\utanyvariationnrm = \utsummationnrm$, then the function considers the sum of all events of activity instances for outcome unit calculations, i.e., it uses $\indu\utsummation$. \\
Overall Quality (case granularity) & \dquality & \gcase & \(c \in C\) & \(\indoq(c)\) & \(\attribute{\atqual}(c)\) \\
Expected Overall Quality & \dquality & \ggroup & \(g \subseteq C \colon\) \newline \(g \neq \emptyset\) & \(\utexpected\indoq(g)\) & \(\frac{\sum _{c \in g} \indoq(c)}{\indc(g)}\) \\
Repeatability (case granularity) & \dquality & \gcase & \(c \in C\) & \(\indrep(c)\) & \(1 - \frac{\inda(c)}{\indi(c)}\) \\
Repeatability (group of cases granularity) & \dquality & \ggroup & \(g \subseteq C \colon\) \newline \(g \neq \emptyset\) & \(\indrep(g)\) & \(1 - \frac{\inda(g)}{\indi(g)}\) \\
Expected Repeatability & \dquality & \ggroup & \(g \subseteq C \colon\) \newline \(g \neq \emptyset\) & \(\utexpected\indrep(g)\) & \(1 - \frac{\sum _{c \in g} \inda(c)}{\indi(g)}\) \\
Rework Count (activity granularity) & \dquality & \gact & \(a \in A\) & \(\indrewc(a)\) & \(\sum _{c \in C} \fnmax(0, \fncount(c, a) - 1)\) \\
Rework Count (case granularity) & \dquality & \gcase & \(c \in C\) & \(\indrewc(c)\) & \(\sum _{a \in A} \fnmax(0, \fncount(c, a) - 1)\) \\
Rework Count (group of cases granularity) & \dquality & \ggroup & \(g \subseteq C \colon\) \newline \(g \neq \emptyset\) & \(\indrewc(g)\) & \(\sum _{c \in g} \indrewc(c)\) \\
Expected Rework Count & \dquality & \ggroup & \(g \subseteq C \colon\) \newline \(g \neq \emptyset\) & \(\utexpected\indrewc(g)\) & \(\frac{\indrewc(g)}{\indc(g)}\) \\
Rework Count by Value (activity granularity) & \dquality & \gact & \(a \in A\) \newline \(\utval \in \mathbb{N}\) & \(\indrewcv(a, \utval)\) & \(\sum _{c \in C} \fnmax(0, \fncount(c, a) - \utval)\) \\
Rework Count by Value (case granularity) & \dquality & \gcase & \(c \in C\) \newline \(\utval \in \mathbb{N}\) & \(\indrewcv(c, \utval)\) & \(\sum _{a \in A} \fnmax(0, \fncount(c, a) - \utval)\) \\
Rework Count by Value (group of cases granularity) & \dquality & \ggroup & \(g \subseteq C \colon\) \newline \(g \neq \emptyset\) \newline \(\utval \in \mathbb{N}\) & \(\indrewcv(g, \utval)\) & \(\sum _{c \in g} \indrewcv(c, \utval)\) \\
Expected Rework Count by Value & \dquality & \ggroup & \(g \subseteq C \colon\) \newline \(g \neq \emptyset\) \newline \(\utval \in \mathbb{N}\) & \(\utexpected\indrewcv(g, \utval)\) & \(\frac{\indrewcv(g, \utval)}{\indc(g)}\) \\
Rework of activities Subset (case granularity) & \dquality & \gcase & \(c \in C\) \newline \(\utsubset \subseteq A \colon\) \newline \(\utsubset \neq \emptyset\) & \(\indrewsub(c, \utsubset)\) & \(\sum _{a \in \utsubset} \fnmax(0, \fncount(c, a) - 1)\) \\
Rework of activities Subset (group of cases granularity) & \dquality & \ggroup & \(g \subseteq C \colon\) \newline \(g \neq \emptyset\) \newline \(\utsubset \subseteq A \colon\) \newline \(\utsubset \neq \emptyset\) & \(\indrewsub(g, \utsubset)\) & \(\sum _{c \in g} \indrewsub(c, \utsubset)\) \\
Expected Rework of activities Subset & \dquality & \ggroup & \(g \subseteq C \colon\) \newline \(g \neq \emptyset\) \newline \(\utsubset \subseteq A \colon\) \newline \(\utsubset \neq \emptyset\) & \(\utexpected\indrewsub(g, \utsubset)\) & \(\frac{\indrewsub(g, \utsubset)}{\indc(g)}\) \\
Rework Percentage (activity granularity) & \dquality & \gact & \(a \in A\) & \(\indrewp(a)\) & \(\frac{\indrewc(a)}{\indi(a)}\) \\
Rework Percentage (case granularity) & \dquality & \gcase & \(c \in C\) & \(\indrewp(c)\) & \(\frac{\indrewc(c)}{\indi(c)}\) \\
Rework Percentage (group of cases granularity) & \dquality & \ggroup & \(g \subseteq C \colon\) \newline \(g \neq \emptyset\) & \(\indrewp(g)\) & \(\frac{\indrewc(g)}{\indi(g)}\) \\
Expected Rework Percentage & \dquality & \ggroup & \(g \subseteq C \colon\) \newline \(g \neq \emptyset\) & \(\utexpected\indrewp(g)\) & \(\frac{\sum _{c \in g} \indrewp(c)}{\indc(g)}\) \\
Rework Percentage by Value (activity granularity) & \dquality & \gact & \(a \in A\) \newline \(\utval \in \mathbb{N}\) & \(\indrewpv(a, \utval)\) & \(\frac{\indrewcv(a, \utval)}{\indi(a)}\) \\
Rework Percentage by Value (case granularity) & \dquality & \gcase & \(c \in C\) \newline \(\utval \in \mathbb{N}\) & \(\indrewpv(c, \utval)\) & \(\frac{\indrewcv(c, \utval)}{\indi(c)}\) \\
Rework Percentage by Value (group of cases granularity) & \dquality & \ggroup & \(g \subseteq C \colon\) \newline \(g \neq \emptyset\) \newline \(\utval \in \mathbb{N}\) & \(\indrewpv(g, \utval)\) & \(\frac{\indrewcv(g, \utval)}{\indi(g)}\) \\
Expected Rework Percentage by Value & \dquality & \ggroup & \(g \subseteq C \colon\) \newline \(g \neq \emptyset\) \newline \(\utval \in \mathbb{N}\) & \(\utexpected\indrewpv(g)\) & \(\frac{\sum _{c \in g} \indrewpv(c, \utval)}{\indc(g)}\) \\
Rework Time (activity granularity) & \dquality & \gact & \(a \in A\) & \(\indrt(a)\) & \(\indlt(a) - \sum _{c \in C} \fnfirstinstlt(c, a)\) \\
Rework Time (case granularity) & \dquality & \gcase & \(c \in C\) & \(\indrt(c)\) & \(\indlt(c) - \sum _{a \in A} \fnfirstinstlt(c, a)\) \\
Rework Time (group of cases granularity) & \dquality & \ggroup & \(g \subseteq C \colon\) \newline \(g \neq \emptyset\) & \(\indrt(g)\) & \(\sum _{c \in g} \indrt(c)\) \\
Expected Rework Time & \dquality & \ggroup & \(g \subseteq C \colon\) \newline \(g \neq \emptyset\) & \(\utexpected\indrt(g)\) & \(\frac{\indrt(g)}{\indc(g)}\) \\
Successful outcome Unit Count (activity instance granularity) & \dquality & \ginst & \(i \in I\) & \(\indsuc\utanyvariation(i)\)
\newline \newline \(\utanyvariationnrm \in \{\utsinglenrm, \utsummationnrm\}\) & \(\begin{cases}
\indu\utanyvariation(i) - \attribute{\atunsunit}(\fnicomplete(i)) & \utif \attribute{\atunsunit}(\fnicomplete(i)) \neq \bot \\
\indu\utanyvariation(i) & \utif \attribute{\atunsunit}(\fnicomplete(i)) = \bot \\
\end{cases}\)
\newline \newline The $\utanyvariationnrm$ indicates that the PPI can take multiple forms, in this case, $\utanyvariationnrm \in \{\utsinglenrm, \utsummationnrm\}$. If $\utanyvariationnrm = \utsinglenrm$, then the function considers single events of activity instances for outcome unit calculations, i.e., it uses $\indu\utsingle$; if $\utanyvariationnrm = \utsummationnrm$, then the function considers the sum of all events of activity instances for outcome unit calculations, i.e., it uses $\indu\utsummation$. \\
Successful outcome Unit Count (activity granularity) & \dquality & \gact & \(a \in A\) & \(\indsuc\utanyvariation(a)\)
\newline \newline \(\utanyvariationnrm \in \{\utsinglenrm, \utsummationnrm\}\) & \(\sum _{i \in \fninst(a)} \indsuc\utanyvariation(i)\)
\newline \newline The $\utanyvariationnrm$ indicates that the PPI can take multiple forms, in this case, $\utanyvariationnrm \in \{\utsinglenrm, \utsummationnrm\}$. If $\utanyvariationnrm = \utsinglenrm$, then the function considers single events of activity instances for outcome unit calculations, i.e., it uses $\indsuc\utsingle$; if $\utanyvariationnrm = \utsummationnrm$, then the function considers the sum of all events of activity instances for outcome unit calculations, i.e., it uses $\indsuc\utsummation$. \\
Successful outcome Unit Count (case granularity) & \dquality & \gcase & \(c \in C\) & \(\indsuc\utanyvariation(c)\)
\newline \newline \(\utanyvariationnrm \in \{\utsinglenrm, \utsummationnrm\}\) & \(\sum _{i \in \fninst(c)} \indsuc\utanyvariation(i)\)
\newline \newline The $\utanyvariationnrm$ indicates that the PPI can take multiple forms, in this case, $\utanyvariationnrm \in \{\utsinglenrm, \utsummationnrm\}$. If $\utanyvariationnrm = \utsinglenrm$, then the function considers single events of activity instances for outcome unit calculations, i.e., it uses $\indsuc\utsingle$; if $\utanyvariationnrm = \utsummationnrm$, then the function considers the sum of all events of activity instances for outcome unit calculations, i.e., it uses $\indsuc\utsummation$. \\
Successful outcome Unit Count (group of cases granularity) & \dquality & \ggroup & \(g \subseteq C \colon\) \newline \(g \neq \emptyset\) & \(\indsuc\utanyvariation(g)\)
\newline \newline \(\utanyvariationnrm \in \{\utsinglenrm, \utsummationnrm\}\) & \(\sum _{c \in G} \indsuc\utanyvariation(c)\)
\newline \newline The $\utanyvariationnrm$ indicates that the PPI can take multiple forms, in this case, $\utanyvariationnrm \in \{\utsinglenrm, \utsummationnrm\}$. If $\utanyvariationnrm = \utsinglenrm$, then the function considers single events of activity instances for outcome unit calculations, i.e., it uses $\indsuc\utsingle$; if $\utanyvariationnrm = \utsummationnrm$, then the function considers the sum of all events of activity instances for outcome unit calculations, i.e., it uses $\indsuc\utsummation$. \\
Expected Successful outcome Unit Count & \dquality & \ggroup & \(g \subseteq C \colon\) \newline \(g \neq \emptyset\) & \(\utexpected\indsuc\utanyvariation(g)\)
\newline \newline \(\utanyvariationnrm \in \{\utsinglenrm, \utsummationnrm\}\) & \(\frac{\indsuc\utanyvariation(g)}{\indc(g)}\)
\newline \newline The $\utanyvariationnrm$ indicates that the PPI can take multiple forms, in this case, $\utanyvariationnrm \in \{\utsinglenrm, \utsummationnrm\}$. If $\utanyvariationnrm = \utsinglenrm$, then the function considers single events of activity instances for outcome unit calculations, i.e., it uses $\indsuc\utsingle$; if $\utanyvariationnrm = \utsummationnrm$, then the function considers the sum of all events of activity instances for outcome unit calculations, i.e., it uses $\indsuc\utsummation$. \\
Successful outcome Unit Percentage (activity instance granularity) & \dquality & \ginst & \(i \in I\) & \(\indsup\utanyvariation(i)\)
\newline \newline \(\utanyvariationnrm \in \{\utsinglenrm, \utsummationnrm\}\) & \(\frac{\indsuc\utanyvariation(i)}{\indu\utanyvariation(i)}\)
\newline \newline The $\utanyvariationnrm$ indicates that the PPI can take multiple forms, in this case, $\utanyvariationnrm \in \{\utsinglenrm, \utsummationnrm\}$. If $\utanyvariationnrm = \utsinglenrm$, then the function considers single events of activity instances for outcome unit calculations, i.e., it uses $\indsuc\utsingle$ and $\indu\utsingle$; if $\utanyvariationnrm = \utsummationnrm$, then the function considers the sum of all events of activity instances for outcome unit calculations, i.e., it uses $\indsuc\utsummation$ and $\indu\utsummation$. \\
Successful outcome Unit Percentage (activity granularity) & \dquality & \gact & \(a \in A\) & \(\indsup\utanyvariation(a)\)
\newline \newline \(\utanyvariationnrm \in \{\utsinglenrm, \utsummationnrm\}\) & \(\frac{\indsuc\utanyvariation(a)}{\indu\utanyvariation(a)}\)
\newline \newline The $\utanyvariationnrm$ indicates that the PPI can take multiple forms, in this case, $\utanyvariationnrm \in \{\utsinglenrm, \utsummationnrm\}$. If $\utanyvariationnrm = \utsinglenrm$, then the function considers single events of activity instances for outcome unit calculations, i.e., it uses $\indsuc\utsingle$ and $\indu\utsingle$; if $\utanyvariationnrm = \utsummationnrm$, then the function considers the sum of all events of activity instances for outcome unit calculations, i.e., it uses $\indsuc\utsummation$ and $\indu\utsummation$. \\
Successful outcome Unit Percentage (case granularity) & \dquality & \gcase & \(c \in C\) & \(\indsup\utanyvariation(c)\)
\newline \newline \(\utanyvariationnrm \in \{\utsinglenrm, \utsummationnrm\}\) & \(\frac{\indsuc\utanyvariation(c)}{\indu\utanyvariation(c)}\)
\newline \newline The $\utanyvariationnrm$ indicates that the PPI can take multiple forms, in this case, $\utanyvariationnrm \in \{\utsinglenrm, \utsummationnrm\}$. If $\utanyvariationnrm = \utsinglenrm$, then the function considers single events of activity instances for outcome unit calculations, i.e., it uses $\indsuc\utsingle$ and $\indu\utsingle$; if $\utanyvariationnrm = \utsummationnrm$, then the function considers the sum of all events of activity instances for outcome unit calculations, i.e., it uses $\indsuc\utsummation$ and $\indu\utsummation$. \\
Successful outcome Unit Percentage (group of cases granularity) & \dquality & \ggroup & \(g \subseteq C \colon\) \newline \(g \neq \emptyset\) & \(\indsup\utanyvariation(g)\)
\newline \newline \(\utanyvariationnrm \in \{\utsinglenrm, \utsummationnrm\}\) & \(\frac{\indsuc\utanyvariation(g)}{\indu\utanyvariation(g)}\)
\newline \newline The $\utanyvariationnrm$ indicates that the PPI can take multiple forms, in this case, $\utanyvariationnrm \in \{\utsinglenrm, \utsummationnrm\}$. If $\utanyvariationnrm = \utsinglenrm$, then the function considers single events of activity instances for outcome unit calculations, i.e., it uses $\indsuc\utsingle$ and $\indu\utsingle$; if $\utanyvariationnrm = \utsummationnrm$, then the function considers the sum of all events of activity instances for outcome unit calculations, i.e., it uses $\indsuc\utsummation$ and $\indu\utsummation$. \\
Expected Successful outcome Unit Percentage & \dquality & \ggroup & \(g \subseteq C \colon\) \newline \(g \neq \emptyset\) & \(\utexpected\indsup\utanyvariation(g)\)
\newline \newline \(\utanyvariationnrm \in \{\utsinglenrm, \utsummationnrm\}\) & \(\frac{\sum _{c \in g} \indsup\utanyvariation(c)}{\indc(g)}\)
\newline \newline The $\utanyvariationnrm$ indicates that the PPI can take multiple forms, in this case, $\utanyvariationnrm \in \{\utsinglenrm, \utsummationnrm\}$. If $\utanyvariationnrm = \utsinglenrm$, then the function considers single events of activity instances for outcome unit calculations, i.e., it uses $\indsup\utsingle$; if $\utanyvariationnrm = \utsummationnrm$, then the function considers the sum of all events of activity instances for outcome unit calculations, i.e., it uses $\indsup\utsummation$. \\
Total Cost and Client count Ratio (activity granularity) & \dquality & \gact & \(a \in A\) & \(\indtcclir\utanyvariation(a)\)
\newline \newline \(\utanyvariationnrm \in \{\utsinglenrm, \utsummationnrm\}\) & \(\frac{\indtc\utanyvariation(a)}{\indcli(a)}\)
\newline \newline The $\utanyvariationnrm$ indicates that the PPI can take multiple forms, in this case, $\utanyvariationnrm \in \{\utsinglenrm, \utsummationnrm\}$. If $\utanyvariationnrm = \utsinglenrm$, then the function considers single events of activity instances for cost calculations, i.e., it uses $\indtc\utsingle$; if $\utanyvariationnrm = \utsummationnrm$, then the function considers the sum of all events of activity instances for cost calculations, i.e., it uses $\indtc\utsummation$. \\
Total Cost and Client count Ratio (group of cases granularity) & \dquality & \ggroup & \(g \subseteq C \colon\) \newline \(g \neq \emptyset\) & \(\indtcclir\utanyvariation(g)\)
\newline \newline \(\utanyvariationnrm \in \{\utsinglenrm, \utsummationnrm\}\) & \(\frac{\indtc\utanyvariation(g)}{\indcli(g)}\)
\newline \newline The $\utanyvariationnrm$ indicates that the PPI can take multiple forms, in this case, $\utanyvariationnrm \in \{\utsinglenrm, \utsummationnrm\}$. If $\utanyvariationnrm = \utsinglenrm$, then the function considers single events of activity instances for cost calculations, i.e., it uses $\indtc\utsingle$; if $\utanyvariationnrm = \utsummationnrm$, then the function considers the sum of all events of activity instances for cost calculations, i.e., it uses $\indtc\utsummation$. \\
Expected Total Cost and Client count Ratio & \dquality & \ggroup & \(g \subseteq C \colon\) \newline \(g \neq \emptyset\) & \(\utexpected\indtcclir\utanyvariation(g)\)
\newline \newline \(\utanyvariationnrm \in \{\utsinglenrm, \utsummationnrm\}\) & \(\frac{\indtc\utanyvariation(g)}{\indcli(g)}\)
\newline \newline The $\utanyvariationnrm$ indicates that the PPI can take multiple forms, in this case, $\utanyvariationnrm \in \{\utsinglenrm, \utsummationnrm\}$. If $\utanyvariationnrm = \utsinglenrm$, then the function considers single events of activity instances for cost calculations, i.e., it uses $\indtc\utsingle$; if $\utanyvariationnrm = \utsummationnrm$, then the function considers the sum of all events of activity instances for cost calculations, i.e., it uses $\indtc\utsummation$. \\
Unwanted Activity count (case granularity) & \dquality & \gcase & \(c \in C\) & \(\induac(c)\) & \(|\lsunwl \cap \fnact(c)|\) \\
Unwanted Activity count (group of cases granularity) & \dquality & \ggroup & \(g \subseteq C \colon\) \newline \(g \neq \emptyset\) & \(\induac(g)\) & \(|\lsunwl \cap (\cup_{c \in g} \fnact(c))|\) \\
Expected Unwanted Activity count & \dquality & \ggroup & \(g \subseteq C \colon\) \newline \(g \neq \emptyset\) & \(\induac(g)\) & \(\frac{\sum _{c \in g} \induac(c)}{\indc(g)}\) \\
Unwanted Activity Percentage (case granularity) & \dquality & \gcase & \(c \in C\) & \(\induap(c)\) & \(\frac{\induac(c)}{\inda(c)}\) \\
Unwanted Activity Percentage (group of cases granularity) & \dquality & \ggroup & \(g \subseteq C \colon\) \newline \(g \neq \emptyset\) & \(\induap(g)\) & \(\frac{\induac(g)}{\inda(g)}\) \\
Expected Unwanted Activity Percentage & \dquality & \ggroup & \(g \subseteq C \colon\) \newline \(g \neq \emptyset\) & \(\utexpected\induap(g)\) & \(\frac{\sum _{c \in g} \induap(c)}{\indc(g)}\) \\
Unwanted activity Instance count (case granularity) & \dquality & \gcase & \(c \in C\) & \(\induic(c)\) & \(|\{i \in \fninst(c) \utst \fnact(i) \in \lsunwl\}|\) \\
Unwanted activity Instance count (group of cases granularity) & \dquality & \ggroup & \(g \subseteq C \colon\) \newline \(g \neq \emptyset\) & \(\induic(g)\) & \(\sum _{c \in g} \induic(c)\) \\
Expected Unwanted activity Instance count & \dquality & \ggroup & \(g \subseteq C \colon\) \newline \(g \neq \emptyset\) & \(\utexpected\induic(g)\) & \(\frac{\induic(g)}{\indc(g)}\) \\
Unwanted activity Instance Percentage (case granularity) & \dquality & \gcase & \(c \in C\) & \(\induip(c)\) & \(\frac{\induic(c)}{\indi(c)}\) \\
Unwanted activity Instance Percentage (group of cases granularity) & \dquality & \ggroup & \(g \subseteq C \colon\) \newline \(g \neq \emptyset\) & \(\induip(g)\) & \(\frac{\induic(g)}{\indi(g)}\) \\
Expected Unwanted activity Instance Percentage & \dquality & \ggroup & \(g \subseteq C \colon\) \newline \(g \neq \emptyset\) & \(\utexpected\induip(g)\) & \(\frac{\sum _{c \in g} \induip(c)}{\indc(g)}\) \\
\bottomrule
\label{PPIQualityDefinitions}
\end{longtable}


\begin{longtable}{>{\hspace{0pt}}m{0.14\linewidth}>{\hspace{0pt}}c>{\hspace{0pt}}c>{\hspace{0pt}}m{0.1\linewidth}>{\hspace{0pt}}m{0.19\linewidth}>{\hspace{0pt}}m{0.473\linewidth}}
\toprule
\textbf{Indicator} & \textbf{D.} & \textbf{G.} & \textbf{F. Par.} & \textbf{Function Name} & \textbf{Formal Definition}\\
\midrule \endfirsthead
\toprule
\textbf{Indicator} & \textbf{D.} & \textbf{G.} & \textbf{F. Par.} & \textbf{Function Name} & \textbf{Formal Definition}\\
\midrule \endhead
\bottomrule \endfoot
\captionsetup{font=small}
\caption{Flexibility dimension PPI definitions} \endlastfoot
Activity and Role count Ratio (case granularity) & \dflexibility & \gcase & \(c \in C\) & \(\indaroler(c)\) & \(\frac{\inda(c)}{\indrole(c)}\) \\
Activity and Role count Ratio (group of cases granularity) & \dflexibility & \ggroup & \(g \subseteq C \colon\) \newline \(g \neq \emptyset\) & \(\indaroler(g)\) & \(\frac{\inda(g)}{\indrole(g)}\) \\
Expected Activity and Role count Ratio & \dflexibility & \ggroup & \(g \subseteq C \colon\) \newline \(g \neq \emptyset\) & \(\utexpected\indaroler(g)\) & \(\frac{\sum _{c \in g} \inda(c)}{\sum _{c \in g} \indrole(c)}\) \\
activity Instance and Human Resource count Ratio (activity granularity) & \dflexibility & \gact & \(a \in A\) & \(\indihrr(a)\) & \(\frac{\indi(a)}{\indhr(a)}\) \\
activity Instance and Human Resource count Ratio (case granularity) & \dflexibility & \gcase & \(c \in C\) & \(\indihrr(c)\) & \(\frac{\indi(c)}{\indhr(c)}\) \\
activity Instance and Human Resource count Ratio (group of cases granularity) & \dflexibility & \ggroup & \(g \subseteq C \colon\) \newline \(g \neq \emptyset\) & \(\indihrr(g)\) & \(\frac{\indi(g)}{\indhr(g)}\) \\
Expected activity Instance and Human Resource count Ratio & \dflexibility & \ggroup & \(g \subseteq C \colon\) \newline \(g \neq \emptyset\) & \(\utexpected\indihrr(g)\) & \(\frac{\indi(g)}{\sum _{c \in g} \indhr(c)}\) \\
Client count (activity granularity) & \dflexibility & \gact & \(a \in A\) & \(\indcli(a)\) & \(|\{\attribute{\atclient}(c) \utst c \in C[a \in \fnact(c)]\}|\) \\
Client count (group of cases granularity) & \dflexibility & \ggroup & \(g \subseteq C \colon\) \newline \(g \neq \emptyset\) & \(\indcli(g)\) & \(|\{\attribute{\atclient}(c) \utst c \in g\}|\) \\
Expected Client count & \dflexibility & \ggroup & \(g \subseteq C \colon\) \newline \(g \neq \emptyset\) & \(\utexpected\indcli(g)\) & \(\frac{\indcli(g)}{\indc(g)}\) \\
Directly-Follows Relations and Activity count Ratio (case granularity) & \dflexibility & \gcase & \(c \in C\) & \(\inddfrar(c)\) & \(\frac{\inddfr(c)}{\inda(c)}\) \\
Directly-Follows Relations and Activity count Ratio (group of cases granularity) & \dflexibility & \ggroup & \(g \subseteq C \colon\) \newline \(g \neq \emptyset\) & \(\inddfrar(g)\) & \(\frac{\inddfr(g)}{\inda(g)}\) \\
Expected Directly-Follows Relations and Activity count Ratio & \dflexibility & \ggroup & \(g \subseteq C \colon\) \newline \(g \neq \emptyset\) & \(\utexpected\inddfrar(g)\) & \(\frac{\sum _{c \in g} \inddfr(c)}{\sum _{c \in g} \inda(c)}\) \\
Directly-Follows Relations count (activity granularity) & \dflexibility & \gact & \(a \in A\) & \(\inddfr(a)\) & \(|\cup_{i \in \fninst(a)} \{\fnact(i') \utst i' \in \fnnext(i)\}|\) \\
Directly-Follows Relations count (case granularity) & \dflexibility & \gcase & \(c \in C\) & \(\inddfr(c)\) & \(|\fndfr(c)|\) \\
Directly-Follows Relations count (group of cases granularity) & \dflexibility & \ggroup & \(g \subseteq C \colon\) \newline \(g \neq \emptyset\) & \(\inddfr(g)\) & \(|\cup_{c \in g} \fndfr(c)|\) \\
Expected Directly-Follows Relations count & \dflexibility & \ggroup & \(g \subseteq C \colon\) \newline \(g \neq \emptyset\) & \(\utexpected\inddfr(g)\) & \(\frac{\sum _{c \in g} \inddfr(c)}{\indc(g)}\) \\
Human Resource count (activity granularity) & \dflexibility & \gact & \(a \in A\) & \(\indhr(a)\) & \(|\fnhres(a)|\) \\
Human Resource count (case granularity) & \dflexibility & \gcase & \(c \in C\) & \(\indhr(c)\) & \(|\fnhres(c)|\) \\
Human Resource count (group of cases granularity) & \dflexibility & \ggroup & \(g \subseteq C \colon\) \newline \(g \neq \emptyset\) & \(\indhr(g)\) & \(|\cup_{c \in g} \fnhres(c)|\) \\
Expected Human Resource count & \dflexibility & \ggroup & \(g \subseteq C \colon\) \newline \(g \neq \emptyset\) & \(\utexpected\indhr(g)\) & \(\frac{\sum _{c \in g} \indhr(c)}{\indc(g)}\) \\
Optional Activity count (case granularity) & \dflexibility & \gcase & \(c \in C\) & \(\indopta(c)\) & \(|\{a \in \fnact(c) \utst \exists c' \in C[a \notin \fnact(c')]\}|\) \\
Optional Activity count (group of cases granularity) & \dflexibility & \ggroup & \(g \subseteq C \colon\) \newline \(g \neq \emptyset\) & \(\indopta(g)\) & \(|\{a \in \fnact(c) \utst c \in g \land \exists c' \in C[a \notin \fnact(c')]\}|\) \\
Expected Optional Activity count & \dflexibility & \ggroup & \(g \subseteq C \colon\) \newline \(g \neq \emptyset\) & \(\utexpected\indopta(g)\) & \(\frac{\sum _{c \in g} \indopta(c)}{\indc(g)}\) \\
Optionality (case granularity) & \dflexibility & \gcase & \(c \in C\) & \(\indopt(c)\) & \(\frac{\indopta(c)}{\inda(c)}\) \\
Optionality (group of cases granularity) & \dflexibility & \ggroup & \(g \subseteq C \colon\) \newline \(g \neq \emptyset\) & \(\indopt(g)\) & \(\frac{\indopta(g)}{\inda(g)}\) \\
Expected Optionality & \dflexibility & \ggroup & \(g \subseteq C \colon\) \newline \(g \neq \emptyset\) & \(\utexpected\indopt(g)\) & \(\frac{\sum _{c \in g} \indopta(c)}{\sum _{c \in g} \inda(c)}\) \\
Role and Variant count Ratio (group of cases granularity) & \dflexibility & \ggroup & \(g \subseteq C \colon\) \newline \(g \neq \emptyset\) & \(\indrolevr(g)\) & \(\frac{\indrole(g)}{\indv(g)}\) \\
Role count (case granularity) & \dflexibility & \gcase & \(c \in C\) & \(\indrole(c)\) & \(|\fnrole(c)|\) \\
Role count (group of cases granularity) & \dflexibility & \ggroup & \(g \subseteq C \colon\) \newline \(g \neq \emptyset\) & \(\indrole(g)\) & \(|\cup_{c \in g} \fnrole(c)|\) \\
Expected Role count & \dflexibility & \ggroup & \(g \subseteq C \colon\) \newline \(g \neq \emptyset\) & \(\utexpected\indrole(g)\) & \(\frac{\sum _{c \in g} \indrole(c)}{\indc(g)}\) \\
Variant Case Coverage (case granularity) & \dflexibility & \gcase & \(c \in C\) & \(\indvcc(c)\) & \(\frac{|\{c' \in C \utst \exists tr \in \fntrace(c')[tr \in \fntrace(c)]\}|}{|C|}\) \\
Variant Case Coverage (group of cases granularity) & \dflexibility & \ggroup & \(g \subseteq C \colon\) \newline \(g \neq \emptyset\) & \(\indvcc(g)\) & \(\frac{|\{c \in C \utst \exists tr \in \fntrace(c)[tr \in \fnvariants(g)]\}|}{|C|}\) \\
Variant count (group of cases granularity) & \dflexibility & \ggroup & \(g \subseteq C \colon\) \newline \(g \neq \emptyset\) & \(\indv(g)\) & \(|\fnvariants(g)|\) \\
\bottomrule
\label{PPIFlexibilityDefinitions}
\end{longtable}

\end{document}
